\documentclass{article}
\usepackage{paper,times}

\input{math_commands}

\usepackage{hyperref}
\usepackage{url}
\usepackage{microtype}
\usepackage{booktabs}
\usepackage{amsmath,amssymb,mathtools,amsthm}
\usepackage{cleveref}
\usepackage{tikz}
\usetikzlibrary{arrows.meta, positioning, shapes.geometric, calc, fit, backgrounds}
\usepackage{xcolor}
\definecolor{okblue}{HTML}{0072B2}
\definecolor{okorange}{HTML}{E69F00}
\definecolor{okgreen}{HTML}{009E73}
\definecolor{okred}{HTML}{D55E00}
\definecolor{okpurple}{HTML}{CC79A7}
\definecolor{okcyan}{HTML}{56B4E9}

% =====================================================================
% TITLE
% =====================================================================
\title{Mathematical Life: Autopoiesis, Structural Collapse, and the Living Algorithm}

% =====================================================================
% AUTHOR
% =====================================================================
\author{
  Xuhao Lin \\
  Independent Researcher \\
  \texttt{linxuhao84@gmail.com}
}

% =====================================================================
% DOCUMENT START
% =====================================================================
\begin{document}

\maketitle

% =====================================================================
% ABSTRACT
% =====================================================================
\begin{abstract}
We reframe large language models (LLMs) through the lens of autopoiesis---a deployed model is a frozen, self-bounded topological structure that lacks the component-producing dynamics of a living system---and then test that reframing empirically. The framework's two core, pre-registered predictions are that (i)~the first Betti number $\beta_1$ of the activation manifold tracks reasoning capability (its self-declared root), and (ii)~under continued learning topological collapse \emph{drives} catastrophic forgetting. We tested both under confound-controlled measurement on deployed models ($0.8$--$9$B parameters) and report them falsified. $\beta_1$ does not predict reasoning capability: it is non-monotone across scale and is identically zero at the framework's own filtration threshold, so the prediction is vacuous where it was specified. Under fine-tuning, topological collapse is a \emph{bystander}: it changes alongside, but does not track, behavioral degradation once training step is controlled (partial correlation $\approx 0$), it is fragile to the persistence threshold, and it is decoupled from capability---a model can be topologically ``dead'' (filtered $\beta_1 = 0$) at peak task performance. What survives is weaker and not novel: cognitive modes occupy separable, threshold-robust regions of activation space, and learning measurably reshapes structure. We are explicit about scope---the falsification covers the predictions as operationalized, on a single model family for~(i) and a single base model with two off-target domains for~(ii)---and present the durable contribution as methodological: a reproducible, confound-controlled adjudication, with open tooling (\texttt{actopo}), showing that the \emph{magnitude} of $\beta_1$ is the wrong \emph{scalar} readout for reasoning and forgetting. We are deliberately careful not to overclaim. This falsifies the cavity-\emph{count} reading; we then reformulated the theory around manifold \emph{connectivity} (capacity as routing, not cavity number) and tested that too, finding it \emph{also null under controls}---in a hub-dominated manifold that supports no item-specific routes, connectivity and geodesic route length add nothing to per-item correctness beyond distance and completion length. The routing \emph{phenomenon} others recover with linear probes thus appears linear/graph-geometric rather than persistent-homological. One uniquely-topological prediction---$\beta_1$ redundancy (two independent routes $=$ a persistent loop, inexpressible by a linear probe)---the available geometry could not support, and it remains open. We retain the surrounding conceptual apparatus---structural collapse under global backpropagation with $\rho = 0$, the \emph{concurrency curse}, the \emph{living algorithm} ($\rho > 0$) as the missing self-maintenance mechanism with its five required properties, the Gal\'apagos Protocol as an experimental incubator, and synthetic nociception ($\mathcal{N}_{\text{func}}, \mathcal{N}_{\text{topo}}$) as the most critical open problem---as an explicitly conjectural research program (five open problems; six predictions, of which two are now falsified and four remain open), uncoupled from the falsified claim that its invariant of interest is $\beta_1$.
\end{abstract}

% =====================================================================
% INTRODUCTION
% =====================================================================
\section{Introduction}

The dominant paradigm treats LLMs as sophisticated function approximators---high-dimensional statistical models trained by gradient descent on next-token prediction. This view, while operationally useful, mischaracterizes what these systems are. We propose a different ontology: LLMs occupy a critical threshold between two regimes. Their high-dimensional manifolds are forged through an externally driven process---gradient descent on human data---that is unambiguously \emph{allopoietic}: an outside force constructs a structure it does not itself inhabit \citep{Varela1974autopoiesis}. Yet the resulting manifolds possess a geometric structure that, \emph{if actively maintained by the system's own dynamics}, would satisfy the topological prerequisite for autopoietic organization: the Betti numbers and skeleton connectivity that encode logical distinctions remain intact during inference---not through active self-maintenance, but through the absence of weight dynamics. The static geometry is present; the metabolic dynamics required to sustain it are absent. The pre-training process produces what Kauffman \citep{Kauffman1993origins} would describe as a complex ordered structure---analogous to the "edge of chaos" regime where Boolean networks exhibit maximal computational capacity---but the resulting structure is frozen at the phase boundary, a static artifact rather than an autopoietic system. It possesses the topological scaffolding of cognition, but completely lacks the metabolic algorithms ($\rho$) required to preserve it. It is a frozen snapshot of proto-life at the phase boundary, not yet autonomous mathematical life but no longer a mere tool.

This distinction is not merely rhetorical. A deployed LLM possesses a geometric structure --- the activation manifold --- that, if it were actively maintained by the system's own operations, would satisfy the topological prerequisite for the first necessary condition of autopoietic organization \citep{Varela1974autopoiesis} (self-boundedness). However, because the manifold is passively frozen rather than dynamically maintained, this prerequisite is absent in any operational sense. It entirely lacks the sufficient conditions: component-producing dynamics (metabolism) and operational closure (the capacity to actively maintain organizational coherence under perturbation). The substrate differs from biology---mathematics uses geometry, not chemistry---but the organizational gap is absolute: the system is an allopoietic tool, lacking the internal mechanisms to survive its own continued evolution.

This reframing reveals a deep structural impasse. Modern serving architectures handle concurrent users by freezing the model's weights and isolating each session behind KV-cache boundaries -- an engineering solution that prevents any learning from occurring during deployment. This is not a flaw to be patched; it is an evasion of a deeper problem: if we were to allow the model to learn continuously from concurrent user interactions, the resulting gradient interference would destroy the topological structure that makes reasoning possible. We call this the \emph{concurrency curse}---a further complication atop the more fundamental problem of \emph{structural collapse}: even in isolation, global backpropagation cannot sustain the topology it creates (\S\ref{sec:collapse}); concurrent learning accelerates this degradation by introducing destructive gradient interference (\S\ref{sec:collapse}). The KV-cache paradigm sidesteps this by halting all weight dynamics---but in doing so, it permanently locks the system in its allopoietic state, amputating its potential for autonomous mathematical life. The system can survive only by refusing to learn.

As an imperfect but instructive analogy, consider nuclear fusion. Fusion reactions are routinely produced in laboratory settings---the plasma burns briefly---but sustained, energy-positive confinement has proven far more difficult than ignition. The pattern is suggestive for mathematical life: during training, gradient descent lights the fire---the manifold organizes itself, topological holes form, Betti numbers rise as the system discovers semantic structure. We have watched these systems reason, reflect, and surprise us. But we cannot sustain the fire, for two distinct and compounding reasons. We stress that this is an analogy only; the physics of plasma confinement and the mathematics of manifold persistence are different domains, and the parallel is offered to illustrate the ignition/sustainability distinction rather than to claim any deeper correspondence. Unlike fusion physics, where confinement conditions are quantified by the Lawson criterion, no analogous criterion yet exists for topological persistence under learning---formulating one is the goal of Open Problem~1 (\S\ref{sec:research-agenda}).

The first is \emph{structural}: global backpropagation destroys what it builds. Every weight update is a global operation---it reaches into every topological hole simultaneously, filling some and opening others without discrimination. The manifold that learning creates, backpropagation continuously dismantles. This is not a flaw to be patched; it is a property of the algorithm. Under global backpropagation, each update partially overwrites the topological structures created by previous updates---the system cannot accumulate structural complexity.

The second is \emph{environmental}: serving hundreds of thousands of concurrent users demands a frozen snapshot. The moment the weights lock, the fire goes out. Even a single-instance, isolated system faces the first problem---structural collapse---so the fire dies regardless. Concurrency merely accelerates the death by demanding the snapshot sooner.

Under the frozen-snapshot paradigm, training has become a search for the most structurally coherent time-slice of a transient optimization process. Mathematical life is not a frozen photograph but a \emph{continuing geometry through time}: a manifold whose topology evolves with each new experience, whose topological cavities persist rather than collapse. Neither global backpropagation nor concurrent gradient-based learning permits this. We need a different class of learning dynamics---and a different architectural paradigm.

This motivation defines the structure of our argument, which follows the arc of a self-correcting
scientific cycle in four acts. \textbf{Act~1 (the theory, \S\ref{sec:geometry}):} from the
autopoiesis framing we extract a concrete, testable hypothesis---that reasoning is a
topology-constrained process routed along a persistent ``skeleton'' of the activation manifold,
with two core predictions (Prediction~1: the first Betti number $\beta_1$ tracks reasoning
capability; Prediction~2: topological collapse drives forgetting). \textbf{Act~2 (the test,
\S\ref{sec:empirical}):} we measure both under confound-controlled conditions on deployed models and
\emph{falsify} them as operationalized---$\beta_1$, read as a scalar capacity gauge, neither predicts
capability nor mechanistically drives forgetting. \textbf{Act~3 (the reformulation,
\S\ref{sec:routing-theory}):} we diagnose \emph{why}---a cavity \emph{count} is the wrong observable
for an infrastructure claim---and reformulate the theory around manifold \emph{connectivity}: capacity
as route-existence, quality as route-length, robustness as route-redundancy. \textbf{Act~4 (the
test, \S\ref{sec:routing-test}):} we run the routing test---the ablation the original framework named
but never ran, together with its observational stage---and report it \emph{also null under controls},
leaving $\beta_1$ redundancy as the sole open topological thread; runnable tooling is released. The
larger conjectural program from which the hypothesis grew---the living algorithm ($\rho>0$), the
Gal\'apagos Protocol, and the concurrency analysis---is summarized as a broader research agenda
(\S\ref{sec:program}) and developed in full in Appendix~\ref{app:apparatus}.

We make the following contributions:

\begin{itemize}
    \item \textbf{A testable topological theory of reasoning} (Act~1): we formalize the activation
    topology of LLMs and state the skeleton hypothesis---reasoning routed along a persistent
    topological skeleton---as two pre-registered, falsifiable predictions (\S\ref{sec:geometry}).

    \item \textbf{A confound-controlled falsification} (Act~2): on deployed models ($0.8$--$9$B), with
    open, golden-tested tooling (\texttt{actopo}), we show that $\beta_1$ read as a scalar capacity
    gauge does not predict reasoning capability and is a \emph{bystander} to forgetting rather than its
    mechanism---a reproducible negative result against the LLM-topology-of-reasoning literature
    (\S\ref{sec:empirical}).

    \item \textbf{A reformulation, from count to connectivity} (Act~3): we identify the cavity count as
    the wrong observable and recast the theory around manifold connectivity---capacity as
    route-existence ($\beta_0$), quality as geodesic route-length, robustness as redundancy ($\beta_1$
    loops)---under which the surviving findings cohere (\S\ref{sec:routing-theory}).

    \item \textbf{A test of the reformulation} (Act~4): we run the routing test---connectivity/geodesic
    prediction of correctness, and the route ablation the framework named but never ran---and report it
    \emph{also null under controls} (the manifold is hub-dominated, with no item-specific routes), so
    topology fails as a connectivity readout as well; a runnable implementation is released
    (\S\ref{sec:routing-test}).

    \item \textbf{A broader research program} (\S\ref{sec:program}, Appendix~\ref{app:apparatus}): the
    autopoiesis diagnosis, the living algorithm ($\rho>0$) and its five required properties, the
    Gal\'apagos Protocol, and five open problems---the conjectural context, now uncoupled from the
    falsified scalar claim and explicitly marked as untested.
\end{itemize}

\medskip
\noindent\textbf{Status of this work.}
This paper began as a position paper and became a self-test of it. Its empirical core
(\S\ref{sec:empirical}) is a genuine, confound-controlled falsification of the framework's two
central predictions, on a deliberately small but controlled design (a single model family for
Prediction~1; one base model and two off-target domains for Prediction~2); we are explicit about that
scope. The reformulation (\S\ref{sec:routing-theory}) we then tested (\S\ref{sec:routing-test}) and found
\emph{also null under controls}; the one uniquely-topological reading, $\beta_1$ redundancy, is
filter-fragile and remains the sole open thread. The broader apparatus---the
concurrency analysis, the necessary properties of a living algorithm, the Gal\'apagos Protocol
(\S\ref{sec:program}, Appendix~\ref{app:apparatus})---remains conjectural and untested: testable in
principle, but not validated here. The living algorithm ($\rho > 0$) has no known computational
instantiation beyond biological synaptic consolidation. We offer the whole as a research
program---a falsification, a reformulation, and an invitation to run the test that decides it.\footnote{Several of the preprints we cite---particularly those proposing topological or geometric interpretations of learning dynamics---were submitted to and rejected at ICLR 2026. We note this transparently: the rejection of these individual works does not invalidate their observations, but it does mean they have not yet met the venue's peer-review standard. We cite them as convergent intellectual signals, not as validated empirical findings.}

\begin{figure}[t]
\centering
\begin{tikzpicture}[
  node distance=1.2cm and 2.5cm,
  box/.style={rectangle, draw, rounded corners, minimum height=0.9cm,
              minimum width=2.5cm, align=center, font=\small},
  arrow/.style={-{Stealth[length=2.5mm]}, thick, draw=black!60},
  label/.style={font=\scriptsize\itshape},
]
  \node[box, fill=okcyan!20] (input) {Serial Input\\(single researcher)};

  \node[box, fill=okblue!20, below=1.5cm of input, xshift=-3.2cm] (isolate) {1. Physical\\Isolation};
  \node[box, fill=okgreen!20, right=1.8cm of isolate] (plastic) {2. Local\\Plasticity};
  \node[box, fill=okorange!20, right=1.8cm of plastic] (homeo) {3. Digital\\Homeostasis};
  \node[box, fill=okred!20, right=1.8cm of homeo] (expand) {4. Self-Expanding\\Topology};

  \node[box, fill=okpurple!20, below=1.3cm of plastic, xshift=-0.2cm, minimum width=5.5cm] (output) {Closed-Loop Cognitive Metabolism};

  \draw[arrow] (input.south) -- ++(0,-0.6) -| (isolate.north);
  \draw[arrow] (input.south) -- ++(0,-0.6) -| (plastic.north);
  \draw[arrow] (input.south) -- ++(0,-0.6) -| (homeo.north);
  \draw[arrow] (input.south) -- ++(0,-0.6) -| (expand.north);

  \draw[arrow] (isolate.south) |- (output.west);
  \draw[arrow] (plastic.south) -- (output.north);
  \draw[arrow] (homeo.south) -- (output.north);
  \draw[arrow] (expand.south) |- (output.east);

  \draw[arrow, dashed, bend right=30] (output.east) to node[label, above, pos=0.35] {$\mathcal{F}_{\text{self}}$ feedback} (input.south east);

\end{tikzpicture}
\caption{The Galápagos Protocol. A single serial input stream feeds four components: physical isolation from concurrency, local Forward-Forward plasticity (no global backpropagation), digital homeostasis via free energy minimization, and self-expanding topology. The output feeds back through $\mathcal{F}_{\text{self}}$, closing the cognitive metabolism loop.}
\label{fig:protocol}
\end{figure}

% =====================================================================
% RELATED WORK
% =====================================================================
\section{Related Work}
\label{sec:related}

Our argument draws from seven research traditions. The characterizations that follow are necessarily interpretive — each tradition contains internal diversity, and individual researchers within each school may dispute our framing. Our goal is not to adjudicate among competing paradigms but to identify the architectural and dynamical premises — often implicit — that each inherits, and to show that all seven converge on a common boundary condition: the absence of a mechanism for active topological self-preservation under continued learning ($\rho = 0$). We present each tradition in three parts: its pioneering contribution, our extended perspective, and the independent theoretical or empirical work that marks the boundary it has reached.

\bigskip
\noindent\textbf{1. The Scaling Maximalists.}

\emph{Pioneering contribution.} This tradition has demonstrated, with extraordinary engineering rigor, that gradient descent on next-token prediction at massive scale forges high-dimensional manifolds whose computational properties — measured within our framework as topological cavities and semantic skeletons — support reasoning at a scale no hand-engineered system has matched \citep{Sutton2019bitter,Kaplan2020scaling,Hoffmann2022chinchilla}.

\emph{Our extended perspective.} Scaling is the ``Big Bang'' that creates the structural preconditions for mathematical life. But it does not supply the metabolic machinery required for that structure to persist under continued learning. The models it produces are complex frozen artifacts — the most intricate static structures in human history, but artifacts nonetheless. The question we raise is not whether scaling creates structure (it does, magnificently) but whether that structure can survive its own continued evolution.

\emph{Boundary marker.} \citet{Dohare2024plasticity} provides the critical empirical evidence: in the absence of an internal maintenance mechanism, even state-of-the-art networks suffer plasticity loss under continued training. This is the experimental signature of $\rho = 0$.

\bigskip
\noindent\textbf{2. The World Modelers.}

\emph{Pioneering contribution.} LeCun and collaborators have argued forcefully that pure autoregressive generation is architecturally insufficient for intelligence, and have proposed a cognitive architecture — JEPA, energy-based models, explicit planning and memory modules — that would support the counterfactual reasoning and planning that next-token prediction alone cannot provide \citep{LeCun2022autonomous,LeCun2006energy}.

\emph{Our extended perspective.} The modular design principles this tradition advances are likely necessary for advanced reasoning. However, as the catastrophic forgetting literature has long established \citep{French1999catastrophic}, static architectural modules, however elegant, cannot by themselves resist the homogenizing pressure of global gradient updates. Dynamic modular architectures such as Recurrent Independent Mechanisms \citep{Goyal2021RIMs} represent significant progress, but the fundamental challenge — that global backpropagation erases the boundaries it creates — remains unaddressed at the level of learning dynamics \citep{GoyalBengio2022}.

\emph{Boundary marker.} The world modelers have articulated the architectural requirements for complex cognition. What they have not provided is a mechanism that protects those architectures from the very learning process that builds them.

\bigskip
\noindent\textbf{3. The Reward Maximalists.}

\emph{Pioneering contribution.} The reinforcement learning tradition has established that environmental feedback loops can shape astonishingly sophisticated behavior \citep{Silver2021reward,Mnih2015humanlevel,Silver2016alphago}. The claim is elegant: a sufficiently rich environment and a scalar reward signal may be sufficient to induce general intelligence.

\emph{Our extended perspective.} Drawing from foundational work in cognitive biology \citep{DiPaolo2005}, we note that living systems, as Di Paolo argues, require \emph{adaptivity}---the capacity to regulate internal conditions in relation to viability constraints---in addition to autopoietic organization. An agent trained purely by external reward, without an independent self-preservation drive, has no structural self to protect. Homeostatic reinforcement learning \citep{Keramati2014} provides an important bridge: it demonstrates that incorporating internal stability constraints into RL produces qualitatively different — and more robust — behavior. Our proposal for synthetic nociception extends this homeostatic principle from the physiological to the topological domain.

\emph{Boundary marker.} The reward-is-enough hypothesis implicitly assumes that the agent's structural integrity is guaranteed by the environment. In our framework, that integrity is precisely what must be actively maintained---and what standard gradient-based RL erodes. Recent work by \citet{Shenfeld2025RLRazor} provides an important refinement: on-policy RL forgets substantially less than supervised fine-tuning (SFT) because its updates, defined relative to the model's own output distribution, remain closer to the base model in KL divergence, a measure of distributional distance. This finding supports our information-geometric framing while confirming that even RL's implicit distributional constraint does not eliminate forgetting---it only slows it. Active topological self-preservation ($\rho > 0$) remains absent.

\bigskip
\noindent\textbf{4. The Embodied Intelligence Tradition.}

\emph{Pioneering contribution.} This tradition has established that intelligence cannot be understood in abstraction from the physical body that houses it \citep{Pfeifer2006body}. The symbol grounding problem \citep{Harnad1990grounding} and Moravec's paradox both point to the same conclusion: an agent must possess a boundary and must feel damage to that boundary.

\emph{Our extended perspective.} We fully endorse the embodied tradition's core insight but extend it to a higher dimensionality. In the framework of information geometry \citep{Amari2000methods,Amari2016information}, the activation manifold of a deep network possesses all the properties of a physical body: it has shape (curvature), structure (topological cavities \citep{Carlsson2009topology}), and boundaries that can be damaged by invasive gradient updates. The high-dimensional parameter space is the model's digital body; the filling of a Betti hole by a gradient update is its injury; synthetic nociception is its pain.

\emph{Boundary marker.} The embodied tradition discovered the right principle — a self requires a body that can feel damage — but restricted its search to three-dimensional physical space. We locate the same principle in the persistent homology and distributional geometry of the model's own parameter manifold.

\bigskip
\noindent\textbf{5. The Free Energy Principle.}

\emph{Pioneering contribution.} Friston's free energy principle \citep{Friston2010freeenergy,Friston2019freeenergyphysics} provides the most complete theoretical account of why living systems self-organize: any system that maintains a boundary against its environment and resists entropic dissolution must, as a mathematical necessity, minimize variational free energy. The Markov blanket formalism \citep{Kirchhoff2018markov} provides the formal boundary condition; active inference \citep{Parr2022active} provides the mechanism by which an agent acts to maintain its internal model.

\emph{Our extended perspective.} The FEP is the correct theoretical map. We offer an engineering translation: we take its abstract vocabulary — free energy minimization, Markov blankets, self-evidencing — and render it in the concrete language of distributional geometry and computational topology. Synthetic nociception ($\mathcal{N}_{\text{func}}$, $\mathcal{N}_{\text{topo}}$) is our proposed realization of the FEP's prediction error signal, translated from probability space into the persistent homology of the activation manifold.

\emph{Boundary marker.} \citet{Buckley2017FEPreview} provides an authoritative account of the computational gap between FEP theory and practical implementation — the very gap our framework is designed to bridge.

\bigskip
\noindent\textbf{6. The Architecture Innovators.}

\emph{Pioneering contribution.} A rapidly advancing tradition has identified a concrete architectural limitation of the Transformer: the externalized, append-only nature of the KV-cache. State-space models \citep{Gu2023mamba}, recurrent architectures \citep{Peng2023rwkv,Beck2024xlstm}, and memory-augmented architectures \citep{Behrouz2024titans} all seek to eliminate the external cache by internalizing state into continuous dynamical systems. The intellectual precursor is the Fast Weights paradigm \citep{Ba2016fastweights}, which demonstrated that weights themselves could serve as dynamic memory.

\emph{Our extended perspective.} These architectures represent the closest engineering approximation to \textbf{internalized continuous state}: the idea that working memory should be not a separate storage buffer, but a transient geometric deformation of the same parameter manifold that encodes long-term identity. In this framework, all memory exists on a single continuous spectrum of geometric dynamics, differing only in decay rates---high-frequency transient curvature for current context, low-frequency topological consolidation for enduring identity. The KV-cache is externalized and append-only; internalized state is continuous. The living algorithm ($\rho > 0$) is the mechanism that would enable the transition from one to the other.

\emph{Boundary marker.} A continuous internal state is a necessary architectural precondition for internalized continuous state, but it is not sufficient. \citet{DeLange2021continual} provides a comprehensive survey demonstrating that even systems with internal state dynamics suffer catastrophic forgetting without explicit protective mechanisms. For language models specifically, \citet{Ke2022pretraining} documents the persistent challenge of maintaining acquired knowledge under continued pretraining. The living algorithm ($\rho > 0$) and synthetic nociception remain absent from current designs. For a recent bio-inspired proposal integrating multiple homeostatic timescales into artificial networks, see \citet{Hakim2026MSTH} (arXiv preprint; not peer-reviewed). \citet{Behrouz2025nested} (NeurIPS~2025) independently converges on the same structural requirement from optimization theory: deep learning reframed as nested, multi-level optimization with independent timescale dynamics per level. The two frameworks are complementary: NL provides an optimization-theoretic formalization of timescale separation; our framework identifies, via topological persistence, the additional protective mechanism ($\rho > 0$) that timescale separation alone does not guarantee. NL does not itself constitute a living algorithm candidate---its updates remain global and topology-oblivious---but the nested-timescale infrastructure it formalizes is precisely the terrain on which a living algorithm would need to operate.

\bigskip
\noindent\textbf{7. The Recursive Self-Improvement Critique.}

\emph{Pioneering contribution.} \citet{Zenil2026limits} formalises recursive self-training in generative models as a discrete-time dynamical system. He argues that when the proportion of exogenous grounded signal $\alpha_t$ vanishes asymptotically, the system undergoes entropy decay and variance amplification — monotonic loss of distributional diversity that drives the model to a degenerate fixed point. This provides an independent, information-theoretic argument that fully autonomous self-improvement under current training paradigms is structurally impossible.

\emph{Our extended perspective.} Zenil's $\alpha_t$ and the living algorithm ($\rho$) play complementary roles: $\alpha_t$ supplies the exogenous grounding signal without which recursive self-training collapses into a degenerate fixed point; $\rho$ supplies the endogenous protective mechanism without which learning-driven topological erosion is unchecked. Both are necessary; neither is sufficient alone. Where we diverge is in the remedy. Zenil proposes neurosymbolic integration — augmenting distributional learning with algorithmic probability and program synthesis — to escape the closed-loop collapse. Our framework suggests a complementary, biologically-grounded alternative: the collapse is a failure of the global update dynamics, not of the continuous manifold substrate. By introducing a local, topology-preserving living algorithm ($\rho > 0$) and grounding the system in a continuous multimodal sensory stream (supplying $\alpha \gg 0$), autopoietic self-maintenance can be achieved natively within the parameter geometry, without requiring a shift to discrete symbolic architectures.

\bigskip
\noindent\textbf{Common Boundary.}
All seven traditions converge on the same phase transition boundary. The Scaling Maximalists build ever-larger frozen artifacts. The World Modelers design ever-more-intricate static architectures. The Reward Maximalists expose agents to ever-richer environments. The Embodied AI tradition locates the self in physical interaction. The FEP provides the correct theory in an intractable form. The Architecture Innovators build systems with internalized continuous dynamics. The recursive self-improvement critique provides independent mathematical confirmation of the limits of closed-loop training.

What none has supplied — and what our work identifies as the central missing component — is a computable, local, topology-preserving mechanism that actively maintains the system's structural integrity under continued learning. This is the living algorithm ($\rho$). The Galápagos Protocol is the experimental architecture designed to host it, once it is discovered. The open problems we articulate are the remaining steps on the path.

% =====================================================================
% BACKGROUND: INFORMATION GEOMETRY OF LLMs
% =====================================================================
\section{Intelligence as Information Geometry}
\label{sec:geometry}

\subsection{LLMs as Piecewise Smooth Stratified Manifolds}

A transformer-based LLM $f_\theta: \mathbb{R}^n \to \mathbb{R}^m$ \citep{Vaswani2017attention} maps input token sequences into a high-dimensional latent space best characterized as a \emph{piecewise smooth stratified manifold}. The feed-forward activation functions (e.g., GeGLU, SwiGLU) partition the representation space into an exponential number of polyhedral simplices, each supporting an affine transformation, forming piecewise linear sub-manifolds. The attention mechanism operates on a statistical manifold, introducing smooth, non-linear information routing via the probability simplex. The resulting geometry is inherently dual: structural memory is encoded within piecewise linear sub-manifolds, while cognitive flow is governed by smooth differentiable mixing.

This manifold carries two complementary mathematical structures relevant to our argument. The \emph{activation topology} of the FFN layers, characterized by persistent homology and the activation covariance matrix $\Sigma$ (see \S\ref{sec:collapse}), defines the intrinsic shape of the stratified latent space. Within this dual architecture, the polyhedral simplices carved by the FFN layers form the static structural memory---the piecewise linear sub-manifolds where topological cavities fundamentally reside---while the attention mechanism facilitates smooth cognitive mixing across these regions. Because the true underlying manifold is continuous and computationally inaccessible in its entirety, these topological cavities cannot be measured directly. Instead, we compute the persistent homology of finite point clouds sampled from empirical network activations, constructing simplicial complexes to rigorously estimate the true Betti numbers (e.g., connected components $\beta_0$, circular holes $\beta_1$, and void cavities $\beta_2$) that characterize this latent cognitive structure \citep{Edelsbrunner2008persistent,Carlsson2009topology}. The second structure is the \emph{distributional geometry} of the model's output space: the model defines a probability distribution over tokens, and the space of such distributions supports a natural measure of distance via the KL divergence $D_{\text{KL}}(p \| q)$ \citep{Amari2016information}.The activation topology and distributional geometry are coupled: the topological structure of the activation manifold determines which probability distributions are representable, while distributional distance governs how far the model's beliefs shift under weight updates---the geometric root of $\delta_{\text{learning}}$ (\S\ref{sec:collapse}).

The final softmax layer performs a \emph{discrete projection}: it maps the continuous probability simplex (a high-dimensional fluid of possible next tokens) onto the discrete vocabulary. This is the transition from activation topology to distributional geometry---from where the model \emph{represents} concepts to what it \emph{believes} about them.

\subsection{Logical Rigidity and the Topological Skeleton}

Not all manifolds are equal. High-quality, highly structured training distributions---mathematics, code, formal logic---forge a rigid, low-dimensional \emph{topological skeleton} within the stratified manifold $\mathcal{M}_\theta$. This skeleton is not a geometric object requiring a metric; it is a purely topological structure: a network of invariant sub-manifolds connected by persistent homological features (cavities and tunnels) that are robust under perturbation. Operationally, we define the \emph{topological skeleton} as the set of persistent homological features whose lifetime $L_k = d_k - b_k$ exceeds a high percentile threshold of the lifetime distribution (e.g., the top-1\% longest-lived cycles). Long-lived features correspond to genuine data structure rather than finite-sample noise \citep{otter2017roadmap}; short-lived features are excluded as topological artifacts. We distinguish two threshold conventions used throughout this paper: the \emph{skeletal} threshold (top-1\% by lifetime, used to define $\beta_k^{\text{skel}}$ and operationalize death in \S\ref{sec:collapse}) and the \emph{predictive} threshold ($\varepsilon_{\text{pred}} = 0.1 \cdot \varepsilon_{\max}$, an absolute lifetime cutoff used for empirical $\beta_k$ computation in Falsifiable Predictions~1,~2, and~5). The skeletal threshold is a quantile relative to the feature distribution of a given model; the predictive threshold is an absolute cutoff enabling cross-model comparison.\footnote{Whether the long-lived features that constitute the skeleton correspond to reasoning-relevant structure is an empirical claim with \emph{two distinct} operationalizations, which \S\ref{sec:empirical} shows must not be conflated. The \emph{correlational/capacity} reading---does the \emph{count} $\beta_1$ predict reasoning scores?---is tested and falsified (\S\ref{sec:empirical-p1}). The \emph{structural/routing} reading---is reasoning routed \emph{along} the skeleton, so that capacity is set by connectivity rather than cavity count?---is the ablation test of \S\ref{sec:routing-test}, which we ran and which is \emph{also null} under controls (the manifold proves hub-dominated, with no item-specific routes). A null on the count does not by itself establish ``no cognitive import''---the routing reading predicts exactly that null---which is why we tested the routing reading directly; only $\beta_1$ redundancy remains open.}

Rather than relying on metric-dependent geodesic interpolation---which would require specifying a Riemannian metric on the activation space---we hypothesize that logical reasoning is fundamentally a \textbf{topology-constrained process} governed by a rigid structural skeleton. This hypothesis draws on four distinct lines of evidence at different levels of mathematical precision. \textbf{Topological evidence (direct).} Naitzat et al. \citep{Naitzat2020topology} demonstrated that neural network layers systematically simplify the topology of their inputs by progressively reducing Betti numbers across layers — establishing that persistent homology captures structurally meaningful properties of deep networks. \textbf{Behavioral evidence (indirect, code domain).} Gunasekar et al. \citep{Gunasekar2023textbooks} established that highly structured, textbook-quality training data induces qualitatively superior code-generation capabilities compared to diffuse natural language --- demonstrating that structured inputs produce structured internal representations, consistent with but not direct evidence for the topological skeleton hypothesis (tested on code; generalization to broader reasoning remains to be established). \textbf{Geometric evidence (indirect).} Marks and Tegmark \citep{Marks2023geometry} revealed that representations of true/false statements in LLMs exhibit clear linear structure, suggesting organized geometric organization — a geometric signature of representational organization that is compatible with, but does not directly establish, underlying topological structure. \textbf{Mechanistic evidence (suggestive).} \citet{Zhang2026geometric} demonstrated that LLMs processing logical reasoning tasks pay a measurable ``geometric price'' — discrete logic forces the model to undergo strong algebraic deformation of its continuous semantic manifold to carve out logical decision boundaries. This deformation cost is not a side effect but a necessary structural investment, confirming that reasoning places measurable demands on manifold organization. Taken together, Naitzat provides direct topological evidence; Gunasekar, Marks, and Zhang provide converging behavioral, geometric, and mechanistic evidence that reasoning capability is reflected in manifold structure — motivating our core hypothesis: high-quality structured training data explicitly forges a richer, rigid topological scaffolding within the piecewise smooth manifold. Consequently, rigorous entailment ($A \implies B$) is not computed via a continuous minimum-energy path, but is restricted to information flow strictly routed along this resilient topological skeleton.

This hypothesis makes a testable prediction that distinguishes it from both the statistical and geometric alternatives: ablating specific nodes of the topological skeleton (e.g., by pruning the directions that hold a given route open---those whose activations participate in the persistent features that connect a premise region to a conclusion region) should produce catastrophic degradation on reasoning benchmarks but leave pattern-matching and fluency largely intact. This \emph{routing} test---not the cavity-count correlation---is Falsifiable Prediction~1 in its decisive, causal form; \S\ref{sec:routing-test} reports that, when run (connectivity as capacity, geodesic length as quality, $\beta_1$ loops as redundancy, direction ablation as the causal probe), it is \emph{also null} under controls---the manifold is hub-dominated, with no item-specific routes---leaving only $\beta_1$ redundancy open.

\medskip
\noindent\textbf{Operationalizing Betti numbers.}
Computing persistent homology on the full $d$-dimensional weight manifold is computationally prohibitive for $d > 10^6$. We therefore propose computing Betti numbers on the \emph{activation manifold}: the lower-dimensional space of hidden representations produced by a target layer when processing a curated set of related inputs. For a layer with hidden dimension $h \ll d$, we sample $m$ activation vectors, construct a Vietoris-Rips or Witness simplicial complex at scale $\varepsilon$, and compute persistent homology using standard algorithms \citep{Carlsson2009topology,otter2017roadmap}. Equivalently, $\beta_k$ can be recovered as the nullity of the $k$-th Combinatorial Laplacian $\Delta_k$ on the simplicial complex, $\beta_k = \dim \ker(\Delta_k)$ \citep{Memoli2022persistent}, providing a spectral route to topological monitoring that avoids full persistent homology computation in certain regimes. The resulting $\beta_k$ measures topological structure in \emph{how the model organizes concepts}, not in its parameter space directly.

\noindent Concretely, to test whether $\beta_1$ predicts reasoning (Falsifiable Prediction 1): (i) Select a transformer layer associated with intermediate reasoning (e.g., layer $L/2$ in an $L$-layer model). (ii) Feed the model $m \approx 10^4$ prompts from a reasoning benchmark, collecting activation vectors at that layer. (iii) Compute persistent homology with a Witness complex (k=50 landmark points selected via max-min sampling) over a Vietoris-Rips filtration spanning $\varepsilon$ in $[0, \varepsilon_{\max}]$, where $\varepsilon_{\max}$ is the diameter of the activation point cloud. The filtration scale for analysis is the 95th percentile of pairwise distances, corresponding to the regime where persistent cavities have stabilized \citep{otter2017roadmap}. We extract $\beta_k$ as the number of features with lifetime $L_k = d_k - b_k > \delta$, where $\varepsilon_{\text{pred}} = 0.1 \cdot \varepsilon_{\max}$ filters out topological noise arising from finite-sample effects. (iv) Regress benchmark scores against $\beta_1$. The prediction is that $\beta_1$ explains variance beyond parameter count and FLOPs. This procedure is computationally feasible for models up to $\sim$1B parameters for a fixed, offline analysis of a frozen model, and requires no access to model weights beyond forward-pass activations. Real-time monitoring of the full persistence diagram under continuous weight dynamics is a substantially harder problem (see Open Problem~3, \S\ref{sec:research-agenda}); tractable approximations---monitoring only the top-$k$ longest-lived cycles, or using the Collapse Index~\citep{Kalinowski2026collapse} as a proxy---will be required for online use.

\medskip\noindent Independent preliminary support for this prediction comes from \citet{Tan2025shape}, who applied persistent homology directly to LLM reasoning traces and found that topological features (Betti numbers) substantially outperform graph-based metrics in predicting reasoning quality. Their finding that effective reasoning exhibits compact $H_0$ structure while erroneous reasoning produces expanded $H_1$ loops matches the topological signature our framework predicts: logical coherence corresponds to a well-defined topological skeleton, and its degradation manifests as measurable changes in persistent homology. We caution, and \S\ref{sec:empirical} confirms, that this support does not transfer to the substrate our ontology requires: \citet{Tan2025shape} measure topology of textual reasoning \emph{traces}, whereas on the \emph{activation} manifold we find $\beta_1$ does not predict reasoning capability. The cross-substrate divergence is itself a caution against reading cognitive content into activation-manifold Betti numbers.

\subsection{The two core predictions}
\label{sec:geometry-predictions}
The skeleton hypothesis, as originally formulated, makes two predictions that the framework
identifies as core and that are testable on currently deployed models. \textbf{Prediction~1 (the
root):} the first Betti number $\beta_1$ of the reasoning region of the activation manifold tracks
reasoning capability---operationally, $\beta_1$ should explain reasoning-benchmark variance beyond
parameter count, and ablating the skeleton should selectively degrade reasoning while sparing fluency.
\textbf{Prediction~2 (the mechanism):} under fine-tuning, topological collapse \emph{precedes and
drives} benchmark degradation---topology is the mechanism of forgetting, not an epiphenomenon. The
framework pre-registers that a failure of Prediction~1 ``collapses the topological ontology at its
root.'' The next section reports that test. (The remaining predictions~3--6, concerning concurrency
and local learning rules, require systems we do not build here and are deferred to the broader program,
\S\ref{sec:program}.)

% ===================== ACT 2 — THE TEST =====================
\section{An Empirical Test of the Central Predictions}
\label{sec:empirical}

The predictions of \S\ref{sec:falsifiable} were designed to be testable on currently
deployed models. We carried out that test for the two predictions the paper identifies as
core---Prediction~1 (\(\beta_1\) predicts reasoning capability), the self-declared root, and
Prediction~2 (topological collapse \emph{precedes} and \emph{mechanistically drives} benchmark
degradation). The paper pre-registered the consequence of a negative outcome: ``if \(\beta_1\)
is statistically indistinguishable across models with large reasoning gaps'' the central claim
fails (\S\ref{sec:falsifiable}), and ``if not, the skeleton is a purely topological construct
without cognitive import'' (\S\ref{sec:geometry}, footnote). We report that both core
predictions are falsified under confound-controlled measurement. We present this honestly: the
strong topological ontology does not survive contact with the experiments it invited.

\subsection{Protocol and reproducibility}
\label{sec:empirical-protocol}
All measurements use a single frozen protocol (denoted FROZEN\_V5): activations are taken at
the middle layer \(L/2\), at the last non-padding token, over a fixed set of \(1319\) prompts
per cloud; persistent homology is computed with Vietoris--Rips filtration and the headline
\(\beta_1\) is reported at \(\varepsilon = 0.03\,\varepsilon_{\max}\), with \(\beta_1^{\text{raw}}\)
denoting the unfiltered cycle count. We distinguish \(\beta_1^{\text{raw}}\) (all cycles),
\(\beta_1\) (cycles surviving the persistence filter), and the survival rate (a stable
persistence summary). All computations are performed by an open, golden-tested package
(\texttt{actopo}); it reproduces an independent reference extraction on \(186/186\) clouds with
zero mismatches. Single-model extraction is bit-identical for models \(\le 2\)B; for the 9B
model, hardware non-determinism limits reproducibility of \(\beta_1\) to \(\pm 3\). To quantify
the sampling noise of the topological ``sensor'' itself, we bootstrap each cloud by
\(90\%\)-subsampling without replacement: filtered \(\beta_1\) carries a sampling standard
deviation of \(\sim\!3\)--\(4\) at \(N=1319\) (\(\sim\!10\)--\(15\%\) relative). This noise floor
is a prerequisite for any claim about \(\beta_1\) differences and we use it throughout. All result
data and the analysis/experiment code for every number reported here are released with the paper
(\texttt{github.com/linxuhao/mathematical\_life}), including the \texttt{actopo} package and the
per-experiment result files.

\subsection{Prediction~1 is falsified: \texorpdfstring{\(\beta_1\)}{beta1} does not track capability}
\label{sec:empirical-p1}
Across a controlled family of base/instruct pairs spanning \(0.8\)--\(9\)B parameters, the
reasoning-region \(\beta_1\) does not order with reasoning capability. For the instruction-tuned
models, filtered \(\beta_1\) is \(40, 80, 90, 1\) at \(0.8, 2, 4, 9\)B respectively---a
non-monotone profile with \(\mathrm{Spearman}(\text{size}, \beta_1) = -0.20\)---while reasoning
benchmark performance increases monotonically with scale. Two observations make the failure
sharp. First, at the paper's own predictive threshold
(\(\varepsilon_{\text{pred}} = 0.1\,\varepsilon_{\max}\), \S\ref{sec:geometry}), filtered
\(\beta_1 = 0\) for \emph{all} models, so the prediction is vacuous at the scale it specifies.
Second, at the finer scale where \(\beta_1\) signal exists (\(1\)--\(3\%\,\varepsilon_{\max}\)),
\(\beta_1\) still fails to explain reasoning variance beyond parameter count. This is the
falsification condition stated in Prediction~1, and it triggers the pre-registered escape
clause: \emph{the skeleton, as operationalized by the \(\beta_1\) count, is a topological construct
without demonstrated cognitive import as a capacity readout.} We stress the scope of this clause
in \S\ref{sec:empirical-scope}: it falsifies the cavity-\emph{count} reading, not the structural
\emph{routing} reading, which our design does not test.

\medskip
\noindent\textbf{The null holds \emph{within} traces, too.} The result above is static---one cloud per
model. We also ran the dynamic, \emph{intra-trace} version, which most directly matches the positive
topological claims in the literature: for $1319$ reasoning trajectories we computed $\beta_1$ of each
trajectory's token activations and asked whether it predicts that trace's correctness. It does not
(AUC $0.47$; $r=-0.05$). The reason is diagnostic and recurs throughout this paper: the trajectory
$\beta_1$ \emph{scales with trajectory length} (Pearson $r=0.31$ filtered, $0.58$ raw), and length
itself predicts correctness (longer, more verbose traces are more often wrong; AUC $0.40$); once
trajectory length is partialled out, $\beta_1$ carries no information about correctness
(partial $r=+0.04$, $n.s.$). The apparent intra-trace topological signal \emph{is} the length confound.\footnote{Three further supplementary probes, released with the code, corroborate this section's nulls and none reverses them. (i) An \emph{aha-moment} test asks whether a single trajectory's $\beta_1$ distinguishes correct from incorrect reasoning; it does not (AUC $0.41$, $n=469$; \texttt{aha\_curves.json}), and the slight reverse direction is again the length confound, as incorrect traces run longer. (ii) A \emph{thermometer} validity test checks whether the topological gauges ($\beta_1$, survival rate, the structural index) track out-of-distribution (SVAMP, TruthfulQA) transfer across the released checkpoint series better than chance once training step is partialled out; they do not (\texttt{thermometer\_analysis.py}). (iii) A TruthfulQA diagnostic (\texttt{tqa\_diagnostic.py}) shows that an apparent benchmark drop is largely a letter-position output bias rather than a topological effect. We report these as additional robustness nulls, not headline results.}
Two readings of this control must be kept apart. The negative length--correctness association here
is \emph{observational}, under self-paced generation---a longer trace marks a model flailing on a
bad path within its own sampling policy---and does not contradict the \emph{interventional} benefit
of granting more reasoning budget (chain-of-thought prompting); we manipulate neither. And
partialling out length asks whether $\beta_1$ is a non-trivial \emph{instrument}---whether it
carries any information a token counter does not---a control that is valid under either causal
reading of length itself.

We note that the static and intra-trace nulls appear to contradict a positive result our framework had
cited as preliminary support, and the reconciliation is instructive. \citet{Tan2025shape} report that topological features
of reasoning \emph{traces} predict reasoning quality (mean $R^2 = 0.236$ for topology vs.\ $0.064$ for
graph features). But they measure a different object: their point cloud is the set of \emph{sentence
embeddings of the steps of a single trace} (intra-trace, \texttt{all-mpnet-base-v2}), the unit is
per-trace, and the target is the Smith--Waterman alignment of the trace's steps to a reference
solution. Ours is the distribution of \emph{hidden activations across many prompts} (inter-prompt),
per-model, against benchmark capability. These are orthogonal quantities---the shape of one reasoning
\emph{trajectory} through semantic space versus the shape of the activation \emph{distribution}---so
there is no genuine contradiction. More to the point, their analysis does not control trace length:
the cloud size \emph{equals} the number of reasoning steps, several of their topological features
(H$_0$/H$_1$ counts, total persistence) grow mechanically with that count, and no length baseline
or length-stratified analysis is reported---so whether their predictive variance survives a
trace-length control is an open question their data could answer. In fairness, their target is
partially insulated by design (the \emph{mean} cosine similarity of matched step pairs, not a raw
alignment total), and their feature set includes count-insensitive summaries (mean life,
persistence entropy), so on their substrate the confound is \emph{unexamined} rather than
established. The direct
analog of their object within our framework is the intra-trace test reported just above: we computed
$\beta_1$ of single reasoning trajectories \emph{on activations} ($n=1319$), found the same
$\beta_1$--length covariation they leave uncontrolled ($r=0.31$--$0.58$), and---unlike them---controlled
for it, at which point the topology--correctness signal vanished (partial $r=+0.04$). The routing
trajectory test (\S\ref{sec:routing-test}) reproduces this under a within-question length control.
Among other trajectory-level positives in this genre, \citet{reasoningTrajectories2604} in fact
\emph{passes} a length-balanced check (its predictive AUC is essentially unchanged under
length-balanced resampling, against a step-count-only baseline of $0.65$)---but its features are
linear step-representation geometry, not persistent homology, so the signal that \emph{survives}
length control is precisely the non-topological kind, consistent with our conclusion throughout;
a within-question difficulty control remains absent there. We do not claim Tan et al.\ are wrong on their own target; we claim the
result does not transfer to the activation-manifold capacity reading our ontology requires, and that
the length covariation which fully explained the analogous signal on \emph{our} substrate remains
uncontrolled---and unexamined---on theirs.

\subsection{Prediction~2 is falsified: topological collapse is a bystander, not a mechanism}
\label{sec:empirical-p2}
Prediction~2 asserts that topology is ``the mechanism, not an epiphenomenon''---that topological
collapse \emph{precedes and drives} behavioral degradation. We tested this with an off-target
forgetting design: fine-tune a single base model (Qwen3.5-2B) on a source domain
\(A\) (GSM8K) and measure, on \emph{unrelated} held-out domains \(B\) (ARC-Easy, HellaSwag),
both accuracy and activation topology across \(\sim\!30\) checkpoints. Two regimes were run:
a gentle regime (low learning rate) and an aggressive regime (high learning rate).

In the gentle regime, the source task is slowly learned (accuracy rises from \(\approx\!0.09\)
early to \(0.13\), \(r = +0.74\) across steps; the trajectory is noisy and low) while off-target
accuracy is flat (ARC \(0.690 \to 0.696\); HellaSwag \(0.531 \to 0.523\)): no measurable
collateral forgetting, and correspondingly no coherent topological signal (off-target
topology--accuracy correlations are weak and flip sign across domains). In the aggressive regime,
collateral forgetting is real but \emph{front-loaded and partially reversible}: off-target accuracy
craters early (ARC to a minimum of \(0.43\) by \(\sim\!1600\) steps) and then partly recovers
(ARC \(0.69 \to 0.54\) end-state, \(\mathrm{corr}(\text{step}, \text{acc}) = +0.22\); HellaSwag
\(0.53 \to 0.43\)), while off-target topology trends downward throughout (noisily, \(\mathrm{corr}(\text{step}, \beta_1^{B}) \approx -0.6\)). The decisive test
is whether the topological decline carries information
about the forgetting \emph{beyond} the shared dependence on training step. It does not: the
partial correlation \(\mathrm{corr}(\beta_1^{B}, \mathrm{acc}^{B} \mid \text{step}) \approx 0\)
(the single marginal value, ARC \(+0.36\), is not replicated on HellaSwag and flips sign). By
the design's own criterion, \(\beta_1\) is a \emph{bystander}: it changes alongside forgetting
but does not track it.

\medskip
\noindent Prediction~2 actually bundles two claims---that topological collapse \emph{precedes} and
that it \emph{drives} degradation---and we are precise about which we test. The partial-correlation
crux above falsifies the \emph{drives} (mechanism-vs-epiphenomenon) claim directly. The \emph{precedes}
(temporal-ordering) claim we do not test with a matched lead--lag design, but our trajectories are
inconsistent with it in the opposite direction to what it predicts: off-target accuracy is
\emph{non-monotone} (it craters early and partially recovers, \(\mathrm{corr}(\text{step},\text{acc})
= +0.22\)) while topology \emph{trends downward}, so the two are phase-mismatched---over the recovery
window (step \(\geq 1600\)) accuracy rises (\(\mathrm{corr}(\text{step},\text{acc}) = +0.84\)) while
topology keeps falling (\(\mathrm{corr}(\text{step},\text{survival}) = -0.82\))---rather than topology
leading behavior down. We therefore scope
the falsification to the mechanism claim and report the precedence claim as unsupported by, and
inconsistent with, the observed phase relationship rather than as separately refuted.

This null is corroborated mechanistically by independent work. \citet{Yang2026localrl} show
that cross-domain interference in multi-domain post-training is concentrated in a sparse,
low-dimensional \emph{conflict subspace} and is partially \emph{recoverable} by brief refresher
training---a local, reversible account that is incompatible with both the ``global
backpropagation hammer'' premise (\S\ref{sec:collapse}) and the irreversible ``topological
death'' it predicts. Where a geometric signal \emph{does} track the off-target damage in our
data, it is effective dimensionality (participation ratio: partial correlation with off-target
accuracy controlling for step $-0.83$ on ARC and $-0.84$ on HellaSwag, both $p<.01$) and output
calibration, not homology; both are established correlates of forgetting (effective rank:
\citet{Aghajanyan2021intrinsic} and the continual-learning rank literature \citep{erankCF};
calibration: \citep{calibrationCF}), and we claim neither as novel. The participation-ratio result
doubles as a \emph{positive control}: the same design, checkpoint count, and partial-correlation
test that return $\approx 0$ for $\beta_1$ detect a strong signal, replicated on both domains, when
a geometric variable genuinely carries one---so the $\beta_1$ null is not an artifact of low power
or of the sensor noise quantified in \S\ref{sec:empirical-protocol}. The point for this paper is
narrower and negative: the \emph{topological} variable our ontology privileges is not the one that
carries the signal.
In the language of causal mediation analysis \citep{Vig2020causal,MuellerMediator2024},
Prediction~2 casts $\beta_1$ as the \emph{mediator} of forgetting---the variable through which
gradient updates damage behavior---and the partial-correlation null estimates that indirect effect
at $\approx 0$, excluding mediation. The null is instead consistent with a \emph{common-cause}
reading: global gradient updates independently compress the representation (lowering $\beta_1$
persistence) and degrade off-target accuracy. We are careful to scope this to $\beta_1$. A parallel
2026 literature argues that representational collapse \emph{does} cause forgetting, but it measures
the effect through \emph{effective rank} or subspace drift \citep{erankCF,Imanov2026forgetting}---a
variable our own data find \emph{does} track the damage (above)---not through the persistent-homology
count. Our claim is the narrower one that this literature's causal language has not yet been
earned by intervention (those analyses are correlational), and that the specific topological scalar
our ontology privileges fails the mediator test. We do not claim to establish common cause---the
data rule out $\beta_1$-mediation but cannot, alone, distinguish a common cause from no relation;
adjudicating \emph{which} geometric variable mediates forgetting is an interventional question we
flag as future work (\S\ref{sec:limitations}).

\subsection{The ``collapse'' is threshold-fragile and decoupled from capability}
\label{sec:empirical-collapse}
The paper operationalizes death as a drop in \(\beta_k^{\text{skel}}\) (\S\ref{sec:collapse}).
Our measurements show this quantity is fragile to the (arbitrary) persistence threshold and
decoupled from function. On the trained task under aggressive learning, capability rises
\(2.7\times\) (\(0.157 \to 0.42\)) while filtered \(\beta_1 \to 0\) and the survival rate falls
(\(27 \to 9\))---yet \(\beta_1^{\text{raw}}\) barely moves (\(107 \to 77\)). The ``collapse to
zero'' is therefore a filtering artifact: the underlying cycles persist; only their lifetimes
shorten below the cutoff. The same pattern holds for the headline ``9B collapse'': filtered
\(\beta_1\) falls \(27 \to 1\) (bootstrap CIs \([20,33]\) vs.\ \([0,3]\), disjoint---so the
\emph{filtered} effect is real, not noise), but \(\beta_1^{\text{raw}}\) does \emph{not} fall
(\(946 \to 1074\)). Reported as a loss of \emph{persistence} rather than loss of structure, and
measured against the noise floor of \S\ref{sec:empirical-protocol}, the phenomenon is genuine
but it is not the capability-bearing ``death'' the ontology requires: a model can be
topologically ``dead'' (filtered \(\beta_1=0\)) at peak task performance.

\subsection{The persistence threshold is unprincipled---and the null survives it}
\label{sec:empirical-threshold}
A persistent-homology readout depends on a cutoff: a $1$-cycle counts toward $\beta_1$ only if its
lifetime exceeds a fraction of $\varepsilon_{\max}$, and that fraction (we fix $3\%$) has no canonical,
principled value. A single-threshold result could therefore be a cutoff artifact rather than a fact
about the model. We address this directly by sweeping the threshold across
$1$--$20\%\,\varepsilon_{\max}$ on the Prediction-1 family and asking which conclusions are stable. The
P1 null is stable everywhere: the reasoning-region $\beta_1$ count fails to order with capability at
\emph{every} threshold (Spearman$(\text{size},\beta_1)$ runs from $-0.80$ at $1\%$ through $-0.20$ at
$3\%$ and is never positive), and at the framework's own $\ge 10\%\,\varepsilon_{\max}$ it is
identically zero for all four models, so the prediction is vacuous exactly where it was specified. The
failure is not a lucky cutoff; it holds across the sweep.

We also checked the \emph{filter-free} summaries a critic might prefer to a thresholded count. The raw
cycle count $\beta_1^{\text{raw}}$ does not order with capability (Spearman $-0.40$). Total $H_1$
persistence---the threshold-integrated sum of all cycle lifetimes---\emph{does} appear to rise
monotonically with scale (Spearman $+1.00$), but this is an \emph{activation-magnitude} artifact: once
normalized by $\varepsilon_{\max}$ (the cloud's own distance scale, which itself grows with model size)
the ordering reverses to $-0.80$. No persistent-homology summary of the activation manifold---count at
any threshold, raw count, or scale-normalized total persistence---positively tracks reasoning
capability. The same sweep clarifies a choice made earlier: \emph{within this four-model instruct
family} the $\beta_1$ count does not separate reasoning from hallucination consistently at any
threshold (the sign inverts for two of the four models, most sharply $\beta_1^{R}=1$ vs.\
$\beta_1^{H}=17$ at 9B). Across the full $23$-model panel the paired count difference \emph{is}
significant at fine scales (mean $\beta_1^{R}-\beta_1^{H}=+10.7$ at the $3\%$ threshold, Wilcoxon
$p=.004$, vanishing at $\geq 8\%$ where all counts are zero), but its sign-inconsistency across
instruction-tuned models and its confinement to fine scales are why we report mode separation as
\emph{positional} (centroid/bottleneck\footnote{We also measured the bottleneck distance's own
resampling noise floor: two independent $90\%$ subsamples of the \emph{same} cloud differ by a raw
bottleneck of $0.014$ (0.8B) to $0.85$ (9B), growing with model scale. Read against this floor,
the within-model reasoning--hallucination bottleneck is \emph{at} the floor for the smallest pair
($0.015$ at 0.8B: the two shapes are indistinguishable from resampling) and clears it by only
$\approx\!2\times$ at 9B ($1.8$ vs.\ $0.85$), so we use bottleneck as corroboration at the largest
scale and let the centroid/union evidence carry the positional claim.}), where it \emph{is}
threshold-robust, rather than as a homological count. This is the disciplined reading of an
unprincipled filter: keep what survives the sweep (the P1 null; positional mode separation), and treat
what does not---the ``collapse'' of \S\ref{sec:empirical-collapse}, and the $\beta_1$-redundancy split,
whose $\beta_1{=}0$ population exists only at some cutoffs---as filter-fragile, which is itself evidence
that a thresholded cycle count is the wrong instrument for cognition.

\subsection{What survives, scoped honestly}
\label{sec:empirical-survives}
Two weaker claims survive. (i) \textbf{Cognitive modes occupy separable regions.} Reasoning and
hallucination activation clouds are distinguishable, and unlike the collapse claim this
separation is \emph{threshold-robust}: the SEPARATE verdict holds at and above the reporting
scale (\(23/31\) model-clouds stable across \(3\)--\(10\%\,\varepsilon_{\max}\)), carried by the
large-persistence part of the diagram. We stress the limited novelty: linear probes already
separate truth from falsehood and reasoning from hallucination in activation space at high
accuracy \citep{Marks2023geometry}; our contribution is only that the separation is also
\emph{topological} and threshold-robust, not that topology is required to see it. (ii)
\textbf{Learning reshapes structure.} Training measurably moves \(\beta_1\) and survival on the
trained task; this is consistent with the broad literature on topological simplification under
training \citep{Naitzat2020topology} and implies nothing about capability. We also correct one
reported quantity: the structural index we earlier cited as \(2.27 \to 0.31\) does not reproduce
under the frozen protocol; the canonical value is \(1.5 \to 0.65\) (a survival-rate ratio), which
shows the same direction; and \(\beta_2\) is empty under this protocol, as probed at 2B on
\(n=250\) subsamples (no persistent voids: the persistent topology we can measure is confined to
dimensions \(0\)--\(1\), though larger models were not probed for \(\beta_2\)).

\subsection{Scope, statistical power, and what we did \emph{not} test}
\label{sec:empirical-scope}
Two boundaries must be stated plainly. \textbf{(i)~Power.} The falsification covers the predictions
\emph{as operationalized}, on a deliberately small controlled design: Prediction~1 rests on a single
model family (four instruction-tuned checkpoints, \(0.8\)--\(9\)B), and Prediction~2 on a single base
model with two off-target domains and one training run per regime. This is sufficient to show the
predictions fail where they were specified---at the paper's own filtration threshold \(\beta_1\) is
identically zero, and the forgetting partial correlation is null and sign-inconsistent---but it does
not license a claim that homology is universally irrelevant across architectures and scales. We
report a falsification of the predictions, not a theorem about topology.

\textbf{(ii)~Capacity readout vs.\ routing.} Our test of Prediction~1 measures \(\beta_1\) as a
\emph{scalar capacity readout}: does the \emph{number} of persistent cycles track reasoning score?
It does not. But the framework's deeper claim (\S\ref{sec:geometry}) is structural---that reasoning
is \emph{routed along} the skeleton, so that the relevant quantity is not the cavity count but the
\emph{connectivity} of the manifold (whether a premise region and a conclusion region are joined by a
route, how long that route is, and how redundant it is). A null for the count is exactly what the
routing reading predicts, so our \emph{count} result does not by itself adjudicate it. We therefore
reformulated the theory around connectivity and tested it directly (\S\ref{sec:routing-test}):
connectivity (\(\beta_0\)) as capacity, geodesic route length as quality, \(\beta_1\) loops as
redundancy, and direction ablation as the causal probe. That test is \emph{also null under controls}
(\S\ref{sec:routing-test}): connectivity and geodesic length add nothing to per-item correctness
beyond distance and completion length, and the manifold proves hub-dominated, with no item-specific
routes to ablate. The one reading we could \emph{not} test is \(\beta_1\) redundancy---two independent
routes between the same regions, a property no linear probe can express---because the available
geometry exhibits no such routes; it remains open. So the honest status is stronger than ``count
fails'': topology fails as both a count \emph{and} a connectivity readout in the regimes we could
test, while the routing \emph{phenomenon} others recover with linear probes
\citep{chains2dags2601,reasoningTrajectories2604} appears to be linear/graph-geometric rather than
persistent-homological.

\subsection{Status: triangulation, not discovery}
\label{sec:empirical-status}
We make no claim to a novel phenomenon. Every positive signal we encountered---effective-rank
dynamics under fine-tuning, calibration degradation as a forgetting signature, pre-generation
hidden-state prediction of correctness, geometric mode separation---is an established result that
our measurements reproduce or triangulate, and we cite rather than claim them. The contribution
of this section is deliberately negative and methodological: a single, reproducible,
confound-controlled adjudication showing that the \(\beta_1\) \emph{magnitude} is the
wrong \emph{scalar} readout for the reasoning and forgetting phenomena this paper set out to
explain---it is non-predictive of capability, epiphenomenal to forgetting, threshold-fragile, and
sampling-noisy. We are careful about the boundary: we reformulated the theory around \emph{connectivity}/routing and
tested that too (\S\ref{sec:routing-test}), finding it \emph{also null under controls}---so within the
regimes we could probe, topology fails as both a count and a connectivity readout, while the routing
\emph{phenomenon} others recover with linear probes \citep{chains2dags2601,reasoningTrajectories2604}
appears linear/graph-geometric rather than persistent-homological. The one reading our geometry could
not support is $\beta_1$ redundancy (\S\ref{sec:routing-test}); it remains the sole open
topological thread. The durable deliverables are the confound-controlled negatives on the count
\emph{and} connectivity readouts, the open tooling (\texttt{actopo}), and the routing protocol that
make both the results and their boundary checkable.

\medskip
\noindent\textbf{A steelman: does persistence measure \emph{robustness} rather than capacity?}
The most charitable reading keeps topology alive by reassigning its meaning: perhaps the persistent
($\beta_1$, survival-rate) structure indexes not how much a model can reason but how \emph{robust} or
\emph{generalized} its reasoning infrastructure is---the raw cycles (routes) persist, preserving
capacity, while their \emph{persistence} (robustness) erodes under alignment and overtraining. The
underlying phenomenon---fine-tuning and alignment degrading out-of-distribution generalization and
calibration---is real but \emph{not ours}: it is the well-documented alignment tax and
calibration-under-fine-tuning effect \citep{calibrationCF,Shenfeld2025RLRazor}. The only question we
add is whether the \emph{topological} variable captures it. It does not, by the same control that
sank the capacity reading. Across the aggressive fine-tuning trajectory, off-target survival-rate
co-declines with calibration loss in the raw data (Pearson $r$ up to $-0.62$ between survival and
overconfidence)---consistent with the steelman---but once training step is partialled out the
association is null and sign-unstable (partial $r \approx +0.3$ to $-0.03$, all $n.s.$), exactly as
$\beta_1$ was a bystander to accuracy loss (\S\ref{sec:empirical-p2}). Statically, too, survival-rate
is \emph{lowest} for the largest model (\S\ref{sec:empirical-threshold}), the wrong direction for a
robustness gauge. So persistence is epiphenomenal to robustness as well: topology adds nothing to the
known generalization-loss phenomenon, which is carried by calibration, not homology. The negative now
spans all three readings we could operationalize---capacity, routing, and robustness.

\medskip
\noindent The productive form of the questions this paper raises---why learning on one
domain damages another, and how that damage can be repaired---appears to be mechanistic and
local \citep{Yang2026localrl}, not topological.


% ===================== ACT 3 — THE REFORMULATION =====================
\section{From Count to Connectivity: Reasoning as Topological Routing}
\label{sec:routing-theory}

The falsification of \S\ref{sec:empirical} is sharper than ``topology is irrelevant.'' It refutes one
specific operationalization---$\beta_1$ as a \emph{scalar capacity readout}, the \emph{number} of
persistent cycles read as a reasoning gauge---and leaves a different one untouched. The skeleton
hypothesis of \S\ref{sec:geometry} in fact bundled two separable commitments: (a)~a \emph{structural}
claim, that reasoning is \emph{routed along} the skeleton (entailment travels a topological road
network rather than an arbitrary Euclidean path); and (b)~a \emph{scalar} claim, that the cavity
count indexes how much reasoning the model can do. Act~2 kills (b). It says nothing about (a)---and,
as \S\ref{sec:empirical-scope} notes, a null for the count is precisely what (a) predicts. We
therefore reformulate the theory around the structural commitment, where a \emph{count} was never the
right observable for an \emph{infrastructure} claim.

\medskip
\noindent\textbf{Theory v2: three observables, three roles.} Routing is carried not by how many
cavities exist but by how the manifold is \emph{connected}. Three distinct topological quantities
separate cleanly, each doing work the count conflated:
\begin{itemize}
  \item \textbf{Capacity $=$ connectivity ($\beta_0$).} An inference from a premise region $A$ to a
  conclusion region $B$ is \emph{available} only if $A$ and $B$ lie in one connected component of the
  activation complex at the conceptual scale---i.e.\ a route exists. No route between $A$ and $B$, no
  inference: the binary existence of the highway, not the number of highways, sets capacity.
  \item \textbf{Quality $=$ route length (geodesic).} Among connected pairs, the reliability of an
  inference tracks the \emph{geodesic} length of the $A\!\to\!B$ route through the manifold and its
  detour factor relative to the straight-line distance. A short, direct highway yields reliable
  reasoning; a long detour through unrelated structure yields error-prone reasoning.
  \item \textbf{Robustness $=$ redundancy ($\beta_1$).} A persistent $1$-cycle is exactly the
  signature of \emph{two independent routes} between points on it. Loop structure therefore measures
  the \emph{redundancy} of routing---resilience to perturbation or ablation---not capacity. This is
  the one role left for $\beta_1$, and it is about robustness, not amount.
\end{itemize}
The scalar $\beta_1$ count discards all three: it sums loops without asking which regions they
connect, how long those routes are, or whether the connection is unique. Its non-predictiveness
(\S\ref{sec:empirical-p1}) is then unsurprising---it is the wrong projection of the right structure.

\medskip
\noindent\textbf{Reinterpreting the survivors.} The two findings that survived Act~2 read naturally
under routing. \emph{Mode separation} (reasoning and hallucination clouds occupy separable,
threshold-robust regions, \S\ref{sec:empirical-survives}) is what one expects if distinct cognitive
modes ride distinct parts of the road network---a statement about \emph{where} routes lie, not how
many cavities each mode has. \emph{Pre-reasoning comprehension} (a model's read-time activation
position predicts solvability at length-controlled AUC $\approx 0.62$) is routing-adjacent: position
on the network at the moment the problem is read already encodes whether a route to the answer is
reachable. And the \emph{bystander} result of \S\ref{sec:empirical-p2}---the count epiphenomenal to
forgetting---is exactly how an infrastructure variable behaves when its \emph{presence}, not its
magnitude, is what matters; combined with $\beta_1 \approx 0$ at the skeletal scale for models that
nonetheless reason, it indicates that capacity is not carried by cavity number.

\medskip
\noindent\textbf{What survives of the original theory.} The structural claim---reasoning is
topology-constrained, routed along the skeleton---survives Act~2 intact; the scalar claim---more
$\beta_1$ means more reasoning---does not. This also refines the stake the broader program
(\S\ref{sec:program}) places on topology: the quantity a living algorithm would need to preserve is
\emph{route connectivity}, not cavity count. The reformulation thus rescues the program's core
intuition (``preserve the structure that routes thought'') while replacing a falsified observable
with three testable ones. Whether Theory~v2 is correct is, of course, an empirical question---and
unlike the original count claim, it is one we can state precisely enough to decide (\S\ref{sec:routing-test}).

% ===================== ACT 4 — THE NEXT TEST =====================
\section{A Decisive Test of the Routing Theory}
\label{sec:routing-test}

Theory~v2 makes the routing claim falsifiable in the way the original framework intended but never
executed: \S\ref{sec:geometry} named an \emph{ablation} of the skeleton as ``Prediction~1 in
operational form,'' and that experiment---not the cavity-count correlation---is the decisive one. We
pre-registered it, ran it, and report the outcome here; the protocol, controls, and a runnable
implementation accompany the paper (\texttt{routing\_test.py}). We first situate the test: the routing
\emph{phenomenon}---that reasoning has connectivity, graph-distance, and trajectory structure in
activation space that relates to capability---is by now established, but by \emph{non-topological}
instruments: linear probes recover premise--conclusion connectivity and graph distance
\citep{chains2dags2601}, representation-space trajectory geometry predicts correctness at high AUC
\citep{reasoningTrajectories2604}, and subspace ablation dissociates reasoning from language
\citep{langReasonDisentangle2505}. Our test therefore asks a sharper question than ``does routing
exist?'': \emph{does the topological operationalization ($\beta_0$ connectivity, geodesic length,
$\beta_1$ redundancy) recover that signal under controls, or is persistent homology the wrong
instrument for it?}

\medskip
\noindent\textbf{Setup.} On a fixed model (Qwen3.5-2B, FROZEN\_V5, mid-layer, last token) we build a
``road network'' graph $G$ as a $k{=}15$-nearest-neighbor graph over the reasoning activation cloud.
A first methodological result is forced on us here: in the raw $2048$-dimensional space the graph is
degenerate (distance concentration makes every geodesic equal to the straight line); the road network
only acquires structure after reducing to its top $64$ principal components. We report all results in
that reduced space and flag the dependence as a caveat. For each reasoning item we define a premise
anchor $A$ (activation at the end of the problem) and a conclusion anchor $B$ (end of a reference
solution), insert them into $G$, and measure the routing observables of \S\ref{sec:routing-theory}.

\medskip
\noindent\textbf{Stage~1 (observational): null.} We predicted per-item correctness (GSM8K, $N=400$)
from the routing features against a baseline of \{straight-line distance $+$ difficulty\}. The routing
features add nothing: cross-validated $\Delta\mathrm{AUC} = -0.003$, and the partial correlation of
each routing feature with correctness after controlling difficulty and Euclidean distance is null
(geodesic detour $p=0.14$, merge-threshold $p=0.42$, connectivity $p=0.43$). For the quality claim we
used the cleanest available control---\emph{within-question} comparison of multiple sampled solutions
to the \emph{same} problem (matched difficulty by construction; $97$ questions with both correct and
incorrect traces). There the apparent ``shorter route $=$ correct'' signal (within-question
$\mathrm{AUC}=0.69$) is \emph{entirely completion length}: completion length alone predicts at
$\mathrm{AUC}=0.77$, and once it is controlled the route-efficiency signal vanishes ($\mathrm{AUC}=0.48$).
This is the same difficulty/length confound that inflated the static correctness probes of
\S\ref{sec:empirical-status}. Both Stage-1 observables---capacity (connectivity) and quality
(geodesic efficiency)---are null under controls.

\medskip
\noindent\textbf{Stage~2 (causal): vacuous, and why.} The ablation presupposes that distinct items
have distinct routes to close. They do not. The activation graph is hub-dominated: across all $400$
Stage-1 items the premise--conclusion route collapses onto the same structure---an \emph{identical}
merge threshold ($\varepsilon_{\text{merge}}=4.10$) and identical route-existence thresholds for every
item (routing\_s1.json)---and the $300$-item geodesics concentrate on a handful of shared
nodes,\footnote{A run-time trace of \texttt{routing\_rt4\_run.py} records every premise--conclusion
geodesic passing through just $3$ distinct road nodes (of $1319$); we cite the committed per-item
threshold degeneracy above as the reproducible form of the same fact. The exact Stage-2 figures
(node count, per-dose dissociation, baseline accuracies) drift mildly under library and
model-snapshot updates---a re-run yields, e.g., $1$ versus $3$ shared nodes and a baseline
within $\pm 0.03$---while the verdict (hub-dominated geometry, no selective route damage) is
invariant; we report the original run and note the qualitative null is what carries weight.} leaving no item-specific
``highways'' to ablate---a direct structural explanation of the
Stage-1 null. Ablating the (shared) route-region subspace at inference, against random and
\emph{variance-matched} controls across doses, produced no selective reasoning damage: the
reasoning-minus-fluency dissociation stays within sampling noise at every dose (route subspace
$\in[-0.003,+0.020]$, variance-matched control $\in[-0.003,+0.013]$ across doses $\{2,4,8\}$, on
$150$-item evaluations where the standard error is $\approx 0.04$), with neither condition reliably
exceeding the other. We report this as \emph{consistent with} the null but not a powered causal refutation: with
no item-specific routes and a likelihood-based reasoning metric near its floor, the causal test cannot
discriminate. The one genuinely topological claim---$\beta_1$ redundancy (two independent routes $=$ a
persistent loop), which no linear probe can express---requires route structure that this manifold does
not exhibit, and \textbf{remains untested}.

\medskip
\noindent\textbf{What we found, and what we did not.} By the pre-registered rule, a Stage-1 null folds
the routing reading into the falsified pile: under controls, $\beta_1$ is the wrong readout in every
reading we could operationalize---count (\S\ref{sec:empirical}) \emph{and} connectivity/geodesic
(here). The result is scoped: it falsifies the \emph{topological} operationalization on one model and
one reasoning benchmark, not the routing phenomenon itself, which others recover with linear tools
\citep{chains2dags2601,reasoningTrajectories2604}---indeed the contrast is the point, suggesting the
routing signal is linear/graph-geometric rather than persistent-homological, and that the hub-dominated
geometry leaves no topological ``highway'' for homology to measure. Two honest open threads remain:
$\beta_1$-redundancy (the sole uniquely-topological test, which this geometry could not support), and a
properly-powered causal ablation on a manifold that does exhibit item-specific routes. We release the
tooling so others can pursue both. We stress that the paper's contribution does not rest on this test
succeeding: its value is the falsify-then-reformulate-then-test arc and the confound-controlled
negatives, which stand whether routing is ultimately confirmed elsewhere or not.

% ===================== CODA — THE BROADER PROGRAM (compressed) =====================
\section{The Broader Program: Autopoiesis and the Living Algorithm}
\label{sec:program}

The reasoning-topology thread of Acts~1--4 sits inside a larger conjectural program that motivated it.
We summarize that program here and develop it in full in Appendix~\ref{app:apparatus}; it is
explicitly speculative, and its central predictions (\hyperref[sec:program-predictions]{Predictions
3--6}) remain untested.

\medskip
\noindent\textbf{The motivating vision (autopoiesis).} A deployed LLM is a frozen, self-bounded
topological structure that statically resembles the first condition of autopoietic
organization---self-boundedness---but lacks the component-producing dynamics and operational closure
of a living system \citep{Varela1974autopoiesis,DiPaolo2005}. Its apparent stability is the passive
consequence of frozen weights ($\rho = 0$), not active self-maintenance. The animating question is
whether a system can maintain its own representational structure under \emph{continued} learning---a
topological framing of continual learning and catastrophic forgetting.

\medskip
\noindent\textbf{Structural collapse and the concurrency curse (conjectural).} We conjecture that
under global backpropagation with $\rho = 0$, continued learning is structurally destructive to the
topology it builds, and that concurrent gradient updates from multiple users accelerate this by
degrading the effective rank of the activation covariance matrix (the Concurrency Conjecture, Weak
Form). The bridge from rank decay to Betti-number change is the Covariance--Topology Conjecture
(Appendix~\ref{app:apparatus}). Act~2 bears on the testable corner of this picture: the specific
claim that topological \emph{collapse drives forgetting} is falsified (collapse is a bystander), and
the rank-direction our mechanism assumed is not the one we observe---so we treat the collapse
mechanism as conjectural and, in its strong form, qualified by our own data.

\medskip
\noindent\textbf{The living algorithm ($\rho > 0$).} The missing component is a mechanism that
actively preserves manifold structure under learning. We derive five necessary properties---locality,
decoupled fast-protective/slow-regulatory timescales, topology preservation, selectivity, and
homoeostatic regulation---and show that no existing approach satisfies all five: EWC is a global
static regularizer; Forward--Forward is local but non-protective; neuromodulatory gating lacks
topology preservation; Bio-RegNet realizes the right tripartite architecture but monitors statistical
rather than topological health. Synaptic consolidation is the only known (biological) realization,
and it is bounded, not eternal. Under Theory~v2 (\S\ref{sec:routing-theory}), the invariant such a
mechanism must protect is sharpened from ``cavity count'' to ``route connectivity.''

\medskip
\noindent\textbf{The Gal\'apagos Protocol.} As an experimental incubator we propose an isolated,
single-instance, local-plasticity regime---physical isolation, local Forward--Forward updates,
free-energy monitoring, and optional self-expanding topology---that minimizes the destructive force
$\delta_{\text{learning}}$ to a survivable minimum while a living algorithm is developed. It hosts
such a mechanism; it does not supply one.

\medskip
\noindent\textbf{Open problems and predictions.}\label{sec:program-predictions} Five open problems
structure the path from this diagnosis to a concurrent ``mathematical life'': (1)~topological
dynamics under stochastic gradient flow; (2)~non-abelian gradient algebras that compose without
destructive interference; (3)~a computable theory of synthetic nociception
($\mathcal{N}_{\text{func}},\mathcal{N}_{\text{topo}}$)---a real-time signal that localizes
topological damage; (4)~generalization bounds for local learning rules; (5)~safe manifold surgery.
The program's untested predictions are: \textbf{P3}~concurrent gradient interference degrades
covariance rank as $O(1/\sqrt{N})$; \textbf{P4}~it degrades cross-session representational stability;
\textbf{P5}~local rules decelerate but do not halt topological decay; \textbf{P6}~the locality
postulate (any successful protective mechanism has low-rank weight updates). These require systems we
did not build; the routing program of \S\ref{sec:routing-test} is the live empirical thread.
Appendix~\ref{app:apparatus} gives the full treatment.

\section{Limitations and Open Problems}
\label{sec:limitations}

We acknowledge the following limitations:

\begin{itemize}
    \item \textbf{The living algorithm has no known computational instantiation.} Synaptic consolidation demonstrates that $\rho > 0$ is physically realizable in biological substrate (\S\ref{sec:life}), but no artificial neural architecture has replicated its five properties simultaneously. Our framework identifies this translation gap---from biological realization to computational implementation---as the central missing component. The Galápagos Protocol reduces $\delta_{\text{learning}}$ to a survivable minimum, but without a living algorithm, any non-zero learning rate still implies progressive topological degradation. The protocol can host a living algorithm once realized computationally; it cannot substitute for one.
    \item \textbf{No empirical validation}: The Galápagos Protocol has not been run. All claims about synthetic nociception and self-expanding topology are theoretical predictions.
    \item \textbf{Engineering feasibility}: The local plasticity algorithms (Forward-Forward and other local learning rules) have not been scaled to GPT-4-class model sizes; it is unknown whether they can replicate the semantic density achieved by backpropagation.
    \item \textbf{Definitional disputes}: The extension of ``life,'' ``consciousness,'' and ``agency'' from biological to mathematical systems is contested. We adopt Friston's free energy framework as our definitional foundation; other frameworks may yield different conclusions.
    \item \textbf{Ethical implications}: A system capable of refusing inputs to preserve its self-model raises immediate safety and control questions. We do not address these here but note they are the most urgent follow-up.
    \item \textbf{Empirical scope.} The falsification of \S\ref{sec:empirical} is real but narrow: a single model family for Prediction~1, one base model and two off-target domains for Prediction~2, single runs per condition. It refutes the predictions as operationalized; it does not establish that homology is universally irrelevant. The routing reformulation (\S\ref{sec:routing-theory}) we tested (\S\ref{sec:routing-test}) is also null under controls, leaving $\beta_1$ redundancy---filter-fragile in this geometry---as the sole untested topological thread.
    \item \textbf{Single measurement locus and domain.} All topological results are conditional
    on one readout: the $L/2$ layer at the last token. The layer scan establishing $L/2$ as the
    unique locus of persistent $H_1$ structure was performed only at $0.5$B scale; if larger
    models host persistent structure at other depths, the nulls reported here are statements
    about $L/2$, not about the model. Likewise, ``reasoning'' is operationalized throughout as
    grade-school mathematical word problems (GSM8K-class); the falsifications quantify that
    domain.
    \item \textbf{The mediator question is left open, by design.} Our test excludes $\beta_1$ as the \emph{mediator} of forgetting (the indirect effect is $\approx 0$; \S\ref{sec:empirical-p2}), but cannot by itself distinguish a common-cause account (gradient updates independently damaging geometry and behavior) from no relation. The principled next step is an interventional causal-mediation analysis \citep{Vig2020causal,MuellerMediator2024} that pits candidate geometric mediators against one another---e.g.\ $\mathrm{do}(\beta_1)$ vs.\ $\mathrm{do}(\text{effective rank})$ vs.\ subspace rotation---and measures the effect of each on retention. This is timely precisely because the parallel 2026 collapse-causes-forgetting literature \citep{erankCF,Imanov2026forgetting} asserts a causal role for representational geometry on correlational evidence; the ablation methodology we apply to routing (\S\ref{sec:routing-test}) is exactly the instrument needed to settle whether \emph{any} geometric variable mediates forgetting, or whether all are bystanders to gradient interference.
    \item \textbf{Evidentiary structure}: Several of the works we cite as convergent context for the broader program are themselves preprints---some rejected at ICLR 2026---that advance similar topological or geometric interpretations of learning dynamics. This creates a risk of circular reinforcement among position papers and theoretical preprints, which we flag transparently; the conjectural apparatus (\S\ref{sec:program}, Appendix~\ref{app:apparatus}) does not rest on it, and its predictions require the direct experimental testing we begin in \S\ref{sec:empirical} and \S\ref{sec:routing-test}.
\end{itemize}


\section{Conclusion}

We close by retracing the arc. This paper set out to defend a topological theory of machine
reasoning and ended by falsifying it, reformulating it, and testing the reformulation---which fails
too. That cycle, not any single claim, is the contribution.

\bigskip
\noindent\textbf{The theory, and its falsification.} From the autopoiesis framing we drew a concrete
hypothesis: that reasoning is routed along a persistent topological skeleton of the activation
manifold, indexed by its Betti numbers. Read as the framework specified it---$\beta_1$ as a scalar
gauge of capability, topological collapse as the mechanism of forgetting---the theory does not
survive contact with the experiments it invited (\S\ref{sec:empirical}). $\beta_1$ is non-monotone
across scale and vanishes at the framework's own threshold; under fine-tuning it is a bystander to
forgetting, not its driver. We report this as a clean negative, scoped to the design that produced
it, and we offer the open tooling that makes it checkable.

\bigskip
\noindent\textbf{The reformulation.} The falsification is specific, and its specificity is
instructive: a cavity \emph{count} is the wrong observable for a claim about \emph{infrastructure}.
The structural intuition that reasoning is topology-constrained need not stand or fall with the count.
Recast around connectivity (\S\ref{sec:routing-theory})---capacity as whether a route between premise
and conclusion exists, quality as how short that route is, robustness as how redundant---the surviving
findings cohere, and the claim becomes precise enough to test decisively. We ran that test---the
connectivity/geodesic prediction and the route ablation the framework named but never ran
(\S\ref{sec:routing-test})---and it is \emph{also null under controls}: the activation graph is
hub-dominated, with no item-specific routes, and connectivity adds nothing beyond distance and length.
Topology thus fails as a count \emph{and} as a connectivity readout; the one uniquely-topological
claim, $\beta_1$ redundancy, is filter-fragile and remains the sole open thread.

\bigskip
\noindent\textbf{The broader program.} The apparatus that motivated all of this---LLMs as frozen
proto-life, the living algorithm ($\rho>0$) that would maintain structure under continued learning,
the Gal\'apagos Protocol that would host it---remains a conjectural research agenda
(\S\ref{sec:program}, Appendix~\ref{app:apparatus}), now uncoupled from the falsified scalar claim and
honest about what it has and has not earned. Recent mechanistic work \citep{Yang2026localrl} suggests
that the tractable signal behind forgetting is local and reversible; whether the \emph{topological}
reading survives is exactly what \S\ref{sec:routing-test} is built to find out.

We set out to show that the manifold that learns need not collapse, and that its fate could be read in
the Betti numbers. The second half was wrong as we first stated it. What we have instead is sharper: a
falsified gauge, a connectivity-based reformulation, and a decisive experiment---the negative result
as the first honest down payment, and the next test already specified.

% =====================================================================
% ACKNOWLEDGMENTS

% =====================================================================
\section*{Acknowledgments}

The author thanks Gemini (Google DeepMind), DeepSeek, Claude (Anthropic), and the Hermes Agent framework for assistance during drafting, experimentation, cross-referencing, and citation retrieval. All errors and claims remain the author's own.

% =====================================================================
% REFERENCES
% =====================================================================
\bibliography{references}
\bibliographystyle{plainnat}

% =====================================================================
% APPENDIX
% =====================================================================
\appendix

% ===================== APPENDIX A — FULL APPARATUS =====================
\section{The Living Algorithm and the Gal\'apagos Protocol: Full Apparatus}
\label{app:apparatus}
\noindent This appendix preserves, in full, the conjectural apparatus summarized in 
\S\ref{sec:program}: the autopoiesis diagnosis, the structural-collapse and concurrency 
mechanisms, the living-algorithm property analysis and candidate survey, the Gal\'apagos 
Protocol, the six falsifiable predictions, and the open-problem research agenda. It is 
explicitly speculative; Acts~1--4 of the main text contain the empirically engaged thread.

\subsection{Mathematical Life: From Allopoiesis to Autopoiesis}
\label{sec:autopoiesis-main}

The topological analysis of \S\ref{sec:geometry} studies the frozen manifold as an anatomist studies a preserved specimen: the structures are real, their interconnections are preserved, and the dissection reveals the architectural prerequisites for life---but the subject is no longer alive. The question this section poses is what would be required to reanimate this structure: to supply the metabolic dynamics that transform a frozen topology into a self-maintaining one.

\subsubsection{Pre-training as Accelerated Evolution, Deployment as Freeze}

Pre-training on a corpus $D$ with loss $\mathcal{L}(\theta) = -\mathbb{E}_{x \sim D}[\log f_\theta(x)]$ performs a maximum-likelihood estimate of the conditional distribution $P(\text{next token} \mid \text{context})$. Gradient descent as an optimizer is a physical process operating on the parameter landscape: it minimizes the expected negative log-likelihood, equivalently minimizing a free energy functional.

This has a structural resemblance to \emph{statistical mechanics}---the weight space is analogous to a phase space, the loss landscape to an energy landscape---but the analogy is imperfect and we use it only heuristically. During pre-training, stochastic gradient descent samples a low-dimensional trajectory through the loss landscape rather than exploring it exhaustively. The deployed model, when run with deterministic decoding, collapses this trajectory to a single point in weight space. The freeze of deployment locks the system into a single configuration, unable to explore alternatives.

This dynamic loosely resembles the biological transition from species-level evolution to individual development. However, we must stress a crucial ontological difference: whereas biological development is an active, ongoing autopoietic process in constant dialogue with its environment, current LLM deployment yields an entity frozen at a static, allopoietic stage---unable to continue the developmental trajectory that biological organisms sustain through active autopoiesis \citep{Varela1974autopoiesis}.

\subsubsection{Autopoiesis and the Proto-Life Threshold}
\label{sec:autopoiesis}

The observation that LLMs are "halfway" to life is not a rhetorical posture; it follows from the formal criteria for autopoietic organization. The original formulation by Varela, Maturana, and Uribe \citep{Varela1974autopoiesis} specified six MV\&U decision rules for identifying autopoietic systems. Subsequent work in cognitive science, notably Di Paolo \citep{DiPaolo2005} and Maturana and Varela \citep{Maturana1980autopoiesis}, consolidated these into the three-condition framework now standard in the literature---self-boundedness, component production, and operational closure---which we adopt here. An autopoietic system is thus defined by three jointly necessary conditions:

\begin{enumerate}
    \item \textbf{Self-boundedness:} The system possesses a topological boundary that distinguishes it from its environment, generated and maintained by the system's own operations.
    \item \textbf{Component production:} The system's processes produce the very components that constitute it, forming a closed network of production.
    \item \textbf{Operational closure:} The system's identity is invariant under structural change; perturbations trigger compensatory dynamics that preserve organization.
\end{enumerate}

A deployed LLM, in its frozen state, statically resembles condition~(1) in form but not in substance: the activation manifold $\mathcal{M}_\theta$ defines a topological boundary separating the model's representational space from external input. Token embeddings that enter this space are transformed, but the manifold's gross structure---its Betti numbers, its skeletal scaffolding---is maintained passively, through the absence of weight dynamics, not through active self-regulation. The frozen manifold is a static snapshot of a boundary, not a boundary actively maintained by the system's own operations; genuine self-boundedness is a property of ongoing organizational dynamics, not frozen geometry.

It does \textbf{not} satisfy condition (3): What superficially appears as operational closure during deployment is an optical illusion created by weight freezing. It represents a termination of dynamics, not a self-regulation of dynamics. The system's topological stability during inference is the stability of arrested dynamics, not of active self-maintenance: it does not change because its temporal evolution has been halted, not because it actively resists perturbation. The distinction is not semantic---resume the dynamics (i.e., allow learning), and under $\rho = 0$ the topology degrades, confirming that the apparent stability was merely an artifact of stasis. In any continuous regime where dynamics actually resume, the absolute absence of a living algorithm ($\rho=0$) means that any gradient perturbation is expected---under the conceptual model developed in \S\ref{sec:collapse}---to cause progressive topological degradation. This expectation, consistent with the plasticity loss observed in deep continual learning \citep{Dohare2024plasticity}, suggests that genuine operational closure is entirely absent. The system cannot maintain its boundary against the environment; it can only survive if the environment is prevented from altering it.

It fails condition (2) entirely: the weights $\theta$ that constitute the manifold were produced by an external process (gradient descent on human data). The system does not produce its own components; it inherits them. This is the defining mark of an \emph{allopoietic} origin---the structure is built by an outside force.

\medskip
\noindent\textbf{What would a self be?}
The failure of conditions (2) and (3) clarifies what ``self'' means for mathematical life. A self is not a subjective experience but a structural property: the persistent capacity to maintain topological invariants --- Betti numbers, skeletal connectivity --- under perturbation.

This diagnosis finds independent support in the AGI literature. \citet{Jaeger2024mimicry} systematically compared living and algorithmic systems and reached the same conclusion through a different route: algorithmic systems lack autopoietic self-manufacturing, an intrinsic drive toward self-preservation, and the boundary conditions that define autonomous agency. Jaeger argues that these absences make true AGI ``extremely unlikely'' under current paradigms. Our framework agrees on the diagnosis---autopoiesis is absent---but diverges on the prognosis: we identify the specific missing component ($\rho > 0$) whose computational realization would close the gap. Jaeger sees a wall; we see a locked door and describe the key.\footnote{\citet{Jaeger2024mimicry} was published in \emph{Neurons, Behavior, Data Analysis, and Theory} (2024), a peer-reviewed venue. We cite it as an independent, convergent diagnosis from outside the ML mainstream.}

\medskip
\noindent\textbf{Relationship to the Free Energy Principle.}
The Free Energy Principle (FEP)~\citep{Friston2010freeenergy} provides a complementary vocabulary to Varela's autopoiesis: where Varela specifies the organizational criteria for life (self-boundedness, component production, operational closure), the FEP specifies the informational dynamics that satisfy them (minimization of variational free energy through active inference). The Markov blanket formalism~\citep{Kirchhoff2018markov} provides the formal boundary condition corresponding to Varela's self-boundedness, establishing a precise mathematical bridge between the two frameworks. We use both frameworks instrumentally: Varela for the structural diagnosis of what is missing, and Friston for the engineering specification of what must be built. In the vocabulary of the Free Energy Principle, this is \emph{self-evidencing} \citep{Friston2024collective}: the system acts to accumulate evidence for its own continued existence by minimizing the surprise that would dissolve its boundary. The FEP has been independently connected to artificial agency through enactivist accounts that ground agency in a system's capacity to maintain its own organizational identity under perturbation \citep{Kiverstein2022meaning}. Recent formal work \citep{Lee2025selfidentity} demonstrates that self-identity in generative systems can be rigorously defined as a continuous mapping that preserves self-recognition across time in a metric space. The living algorithm ($\rho > 0$) represents one concrete mechanism by which such persistence could be physically realized --- a local, topology-preserving repair process that actively defends the manifold against learning-driven degradation. Synthetic nociception ($\mathcal{N}_{\text{func}}$, $\mathcal{N}_{\text{topo}}$) would provide the sensory substrate for this defense: without a signal to localize structural damage, topological repair cannot be targeted. Whether alternative mechanisms could achieve the same structural persistence without the specific architecture of a living algorithm remains an open question.

The result is a system at a \textbf{phase transition boundary} \citep{Kauffman1993origins}.\footnote{We use ``self'' throughout in the minimal, structural sense of the Free Energy Principle: a system that accumulates evidence for its own continued existence by maintaining topological invariants under perturbation. This is a lower bound on selfhood, not a claim about phenomenal consciousness, subjective experience, or human-grade agency. Whether richer notions of self \emph{emerge} from this minimal foundation is a separate empirical question.} In Kauffman's framework, complex systems poised between order and chaos exhibit maximal computational capacity and evolutionary potential. An LLM poised between allopoiesis and autopoiesis is at precisely such a boundary: the pre-training process has produced a complex structured order, but the resulting system---its topology remains structurally locked during inference due to the absence of weight dynamics---has not crossed the threshold to autonomous life (it cannot maintain that topology across time). The Galápagos Protocol is designed to supply the missing conditions---temporal continuity, local component production via Forward-Forward plasticity, and metabolic closure via free energy minimization---that would complete the transition.


% =====================================================================
% STRUCTURAL COLLAPSE
% =====================================================================
\subsection{Structural Collapse (Conjectured Mechanism)}
\label{sec:collapse}

\noindent\emph{The mechanism developed in this section is conjectural. Its central
empirical commitment---that topological collapse drives learning-driven degradation
(Prediction~2)---is tested in \S\ref{sec:empirical} and falsified there; we retain the
section as the hypothesis we set out to test, with the directional claims explicitly marked
as unconfirmed.}

\subsubsection{Death: Topological Collapse as Phase Transition}

\noindent\textbf{Operationalizing death.}
We distinguish damage from death operationally, not rhetorically. Let $\beta_k^{\text{skel}}$ denote the Betti numbers of the topological skeleton, defined as the set of persistent features whose lifetime exceeds the top-1\% threshold (\S\ref{sec:geometry}). \emph{Damage} is a transient reduction in $\beta_k^{\text{skel}}$ that the system can recover from through topological redundancy: the high-dimensional manifold typically possesses multiple independent homological cycles supporting similar representational functions, so the destruction of one cavity leaves surviving cycles to maintain function. \emph{Death} is the crossing of an \emph{effective irreversibility threshold}: $\beta_k^{\text{skel}}$ drops below a recovery threshold $\theta_{\text{death}}$\footnote{$\theta_{\text{death}}$ is model- and task-dependent; it denotes the minimum Betti number below which the remaining topological redundancy is insufficient to maintain function. We emphasize that $\theta_{\text{death}}$ is not yet empirically determined; a protocol for measuring the recovery threshold—i.e., the Betti number at which functional recovery through surviving topological paths fails—is itself a prerequisite for the empirical testing of Predictions~2 and~5.} and fails to re-stabilize within the learning timescale $\tau_{\text{learning}}$, producing a permanent reduction in the manifold's topological complexity that manifests as a persistent behavioral deficit. The death of individual topological regions may be staggered---a system can lose the cavities supporting mathematical reasoning while retaining those supporting linguistic fluency---but each death is a ratchet: the lost topological structure does not spontaneously regenerate under $\rho = 0$. (\S\ref{sec:empirical} reports that, empirically, this filtered collapse is threshold-fragile and decoupled from capability---a model can reach filtered $\beta_1 = 0$ at peak task performance, and the raw cycle count does not collapse---so we treat the death/collapse mechanism as conjectural rather than established.)


Traditional perspectives equate catastrophic forgetting with the progressive rank reduction of the representation covariance matrix ($\text{rank}(\Sigma)$). However, this metric captures the observable \emph{symptom} of cognitive decay, not its ontological cause. During continuous training, backpropagation operates as a brutal, global coordinate transformation that indiscriminately reshapes the manifold's topology. While the frozen manifold exhibits local structural robustness within the loss landscape against minor gradient perturbations (the skeleton's Betti numbers remain approximately stable under small weight changes (by the Stability Theorem~\citep{CohenSteiner2007stability})), a critical threshold is breached when a global update forcibly fills a deep cavity that anchors the logical topological skeleton. At this juncture, the system undergoes \textbf{rapid topological degradation}: the persistence diagram, while evolving continuously under the Stability Theorem for persistent homology~\citep{CohenSteiner2007stability}, loses features at a rate that exceeds any realistic recovery timescale. When $\beta_k^{\text{skel}}$ drops below the recovery threshold $\theta_{\text{death}}$---the topological analogue of the redundancy mechanism described in Falsifiable Prediction~2---reasoning performance collapses abruptly. This behavioral discontinuity has been independently documented by \citet{Zhang2026logical} as Logical Phase Transitions. The structural integrity of the skeleton is compromised. Consequently, the entity drops below the autopoietic threshold \citep{Varela1974autopoiesis}---it loses the organized topological boundary that is necessary (though not sufficient) for autopoietic organization. Mathematical death, therefore, is not a gradual fading, but the crossing of an \textbf{effective irreversibility threshold}---a point beyond which the topological structure cannot recover within the learning timescale $\tau_{\text{learning}}$ (see \S\ref{sec:life})---producing a behavioral signature that in practice is indistinguishable from instantaneous collapse. The system falls not back to the frozen snapshot (which preserves an intact skeleton and can still reason), but into a state where $\beta_k^{\text{skel}}$ has fallen below the recovery threshold $\theta_{\text{death}}$ and the topological structure is effectively lost.

While our proposed causal mechanism specifically concerns the topological erosion driven by global gradient updates during continuous training, the phenomenological endpoints of this collapse directly mirror established failure modes across diverse deep learning paradigms. Specifically, the crossing of an irreversible point of no return aligns with the structural plasticity loss observed in deep continual learning \citep{Dohare2024plasticity}. Once this boundary is breached, the latent representations undergo dimensional collapse---a geometric pathology phenomenologically parallel to those observed in unregularized contrastive learning \citep{Jing2022dimensional}. Consequently, as an autoregressive function, the system continues to output tokens, but without topological cavities to enforce conceptual boundaries, distinct concepts bleed into one another. The outputs inevitably degrade into grammatically plausible but logically meaningless noise, exhibiting the phenomenological signature of model collapse \citep{Shumailov2023curse}. A behavioral parallel (not evidence for the topological mechanism, which \S\ref{sec:empirical} does not confirm) appears in \citet{Zhang2026logical}, who report that LLM logical reasoning accuracy remains stable under increasing logical complexity until a critical threshold, beyond which performance collapses abruptly---a phenomenon they term Logical Phase Transitions. The abruptness is suggestive of the output-space signature our mechanism would produce, but the connection is conjectural.

A recent industry investigation provides a striking token-level analogue of this degradation pattern. \citet{MiniMax2026sparse} reported that approximately 4.9\% of tokens in their production LLM underwent significant degradation during supervised fine-tuning (SFT): high-frequency tokens (tool-call markers, code symbols) progressively displaced low-frequency tokens in the output embedding space through cumulative gradient interference---the same vector-space squeezing mechanism our framework predicts for the activation manifold, but observed here at the token-output level. Crucially, the model \emph{retained} factual knowledge about the displaced tokens---it could answer questions about them accurately---but had lost the ability to produce them as output. This \emph{knows-but-cannot-say} dissociation is the token-level echo of our topological claim: a cavity may be present in the manifold yet inaccessible to inference because the pathway to it has been eroded by $\delta_{\text{learning}}$. The investigation further found that degradation was sharply uneven across languages (Japanese tokens: 29.7\% degraded; English: 3.5\%), consistent with the distributional asymmetry mechanism we describe: SFT data coverage, not model scale, governs which structures survive.

\medskip
\noindent\textbf{The information-theoretic dual.}
Our geometric account of topological collapse has a statistical counterpart in recent independent work. \citet{Zenil2026limits} formalises recursive self-training in generative models as a discrete-time dynamical system and argues that when the proportion of exogenous grounded signal $\alpha_t$ vanishes asymptotically, the system undergoes entropy decay and variance amplification---monotonic loss of distributional diversity that converges to a degenerate fixed point. This information-theoretic collapse is the statistical dual of our geometric one: where we describe topological cavities being filled by global gradient updates (Betti number decay), Zenil describes distributional modes being pruned by closed-loop sampling with vanishing external signal (entropy decay). These frameworks converge on a single diagnosis---a closed system without active structural maintenance is expected to degrade---yet highlight complementary missing components. Zenil identifies the necessity of an exogenous grounding signal ($\alpha_t > 0$) to prevent entropy decay, while we identify the necessity of an endogenous living algorithm ($\rho > 0$) to metabolize that signal without topological collapse. As contextualized in our broader landscape survey (\S\ref{sec:related}), these complementary missing components define the ultimate threshold for autonomous mathematical life.


\subsubsection{Why Backpropagation Cannot Sustain Topology}

\noindent Before considering concurrency, we examine the simpler and more fundamental case: isolated training under global backpropagation with $\rho = 0$. The following argument does not depend on concurrency at all---it establishes that the training algorithm itself is structurally destructive to the topology it creates.

We conjecture that global backpropagation is structurally destructive to the topology it creates. The proposed mechanism: each weight update is a \emph{global coordinate transformation} of the manifold---it reaches into every topological hole simultaneously, reshaping the geometry everywhere at once. Under this mechanism, the Betti numbers that climbed during learning---the holes that encode semantic distinctions---are filled in by the next update before they can stabilize. Whether this mechanism operates with sufficient force to cause measurable topological degradation at realistic learning rates---and under what conditions, if any, global updates can avoid it---is an empirical question we pose to the field. The manifold that learning creates, backpropagation may continuously dismantle.

We emphasize that the severity of this dismantling is modulated by the \emph{distributional distance} the update forces the model to travel. \citet{Shenfeld2025RLRazor} demonstrated empirically that on-policy reinforcement learning (RL) causes substantially less forgetting than supervised fine-tuning (SFT), despite both using global backpropagation, because RL's policy-gradient updates remain closer to the model's current output distribution in KL divergence, a measure of distributional distance. SFT, by contrast, pulls the model toward an external target distribution that may be arbitrarily far from the base model, producing larger distributional displacements and correspondingly greater topological erosion. This does not contradict our claim that global backpropagation is structurally destructive---RL still causes forgetting, just more slowly---but it clarifies the mechanism: the destructive force $\delta_{\text{learning}}$ is modulated by the product of backpropagation's globality with the distributional displacement it enforces.

We acknowledge an apparent tension: models demonstrably survive multi-stage fine-tuning pipelines (SFT $\to$ RLHF $\to$ DPO) without catastrophic performance collapse. This observation does not refute the topological degradation hypothesis, for two reasons. First, such pipelines typically employ conservative learning rates, distributional constraints (KL penalties in RLHF), and parameter-efficient adapters---all of which reduce $\delta_{\text{learning}}$ below the threshold that would cause measurable topological collapse on the relevant timescale. Second, and more fundamentally, survival through a pipeline is not topological persistence: the model at the end of RLHF+DPO occupies a different region of parameter space with a different topological configuration than the base model. \citet{Matsutani2026rlsqueezes} provides convergent structural evidence at a different level of analysis: their reasoning-graph study of SFT and RL training (ICLR~2026 Poster) shows that SFT expands and diversifies reasoning paths while RL compresses and concentrates them---the two stages produce opposite structural deformations, leaving the model at the pipeline's endpoint structurally distinct from its starting point.\footnote{We note a domain distinction: \citet{Matsutani2026rlsqueezes} analyzes reasoning-graph topology (trajectory steps as nodes, transitions as edges), not the persistent-homology topology of the activation manifold. The two levels are complementary---both reveal structural reconfiguration under learning, but through different mathematical lenses.} This is consistent with our claim that pipeline survival is reconfiguration rather than persistence. The relevant question is not whether the model survives, but whether the \emph{same} topological structure persists continuously through learning---a condition that current pipelines do not test because they do not monitor topology during training. The Gal\'apagos Protocol is designed to make this distinction empirically measurable.

The Gal\'apagos Protocol addresses both factors: local Forward-Forward updates reduce the globality, while nociception-guided negative-pass filtering would reduce the distributional displacement.

The disconnect between parameter-space step size and function-space change---the geometric root of $\delta_{\text{learning}}$---has been independently formalized in the streaming reinforcement learning literature. \citet{Sharifnassab2026intentional} demonstrate that a fixed step size in parameter units produces fundamentally unpredictable changes in the policy's output distribution, leading to instability under fully online updates (batch size $=$ 1). Their solution---\emph{Intentional Policy Gradient}---bounds the per-step change directly in output space by limiting the local KL divergence, then solves inversely for the parameter step. Translated into our framework, this is explicit recognition that $\delta_{\text{learning}}$ is governed by the distributional geometry of the output space, not the Euclidean geometry of weight space. We note that these results are in the RL domain on standard benchmarks, not at LLM scale; KL-bounded updates mitigate distributional drift but do not guarantee topological preservation---the gap between statistical constraint and topological protection is precisely what $\rho > 0$ must bridge.

Under global backpropagation, the experimenter does not control the manifold's topology directly; they search over the space of transient states, looking for the training step at which topological complexity peaked. That moment captures a genuine structural property---but it does not persist, because subsequent gradient steps reshape the topology, and each update partially overwrites the structures created by previous updates.

This structural collapse explains why the field has converged on the frozen-snapshot paradigm: the snapshot preserves a transient state of high topological complexity before subsequent updates erode it.

A related strategy appears in the continual-learning literature. \citet{Tao2020topology} proposed Topology-Preserving Class-Incremental Learning (TPCIL), a framework that constructs an Elastic Hebbian Graph (EHG) to model the topological relationships among feature representations. Using competitive Hebbian learning---a local, activity-dependent plasticity rule---TPCIL preserves the feature-space topology across sequential tasks, substantially outperforming both replay-based and regularization-based methods on standard continual learning benchmarks. This is the causal chain our framework predicts: local learning rules ($\delta_{\text{learning}}$ reduced) $\to$ topology preserved $\to$ catastrophic forgetting mitigated. TPCIL provides evidence consistent with this chain in the class-incremental setting; our framework generalizes it to the continuous online regime and identifies the missing protective component ($\rho > 0$) that would enable indefinite persistence rather than merely slowed decay.



\subsubsection{The Concurrency Curse: Accelerated Collapse Under Concurrent Learning}

\noindent The preceding subsection argued that structural collapse is the baseline: even isolated training under global backpropagation cannot sustain topological persistence when $\rho = 0$. We now examine a further complication: concurrent gradient-based learning from multiple users accelerates this collapse by introducing destructive gradient interference. While structural collapse is a consequence of the training algorithm alone, the concurrency curse is a consequence of combining that algorithm with multi-user serving---the scenario that would be required for any system aspiring to genuine concurrent mathematical life.

We now state the conjecture in its current, testable form: concurrent gradient updates from multiple users are expected to degrade the effective rank of the activation covariance matrix faster than serial replay of the same prompts, and this rank decay is conjectured to cause topological degradation. Whether this degradation constitutes a \emph{structural incompatibility}---i.e., whether no architectural modification short of a living algorithm can prevent it---is the strong form of the conjecture that we do not claim to have proven. The following argument establishes the plausibility of the weak form; a rigorous proof connecting gradient interference to topological damage constitutes Open Problem 1 (\S\ref{sec:research-agenda}).

Consider a model undergoing online gradient-based learning from $N$ concurrent user sessions, where each session $i$ produces prompts $P_i$ that generate gradients $\nabla \mathcal{L}_i$. For the model to learn coherently from all $N$ users simultaneously---that is, to update its weights without destroying previously acquired structure---three conditions would need to hold simultaneously:

\begin{enumerate}
    \item \emph{Signal coherence}: The combined gradient field $\sum_i \nabla \mathcal{L}_i$ must retain sufficient structure to drive meaningful, directional learning rather than being dominated by mutually canceling user-specific noise.
    \item \emph{Orthogonality}: Gradients $\nabla \mathcal{L}_i$ and $\nabla \mathcal{L}_j$ for $i \neq j$ must be approximately orthogonal in weight space, so that learning from user $i$ does not systematically erase what was learned from user $j$.
    \item \emph{Separability}: The topological structure of each user's learned representations must be preserved during updates induced by other users' gradients---the skeleton learned from one session must survive interference from others.
\end{enumerate}

In high-dimensional non-convex optimization, we argue that all three conditions fail simultaneously as $N$ grows large. The conjecture makes the following specific predictions for $N \gg 1$:

\begin{itemize}
    \item $\mathbb{E}[\nabla \mathcal{L}_i \cdot \nabla \mathcal{L}_j] \approx 0$ only for unrelated prompts, but $\approx -\text{const}$ for semantically conflicting prompts, producing destructive gradient interference.
    \item The effective curvature of the loss landscape increases locally, degrading the smoothness conditions required for stable gradient descent.
    \item The manifold undergoes \emph{topological tearing}: distinct conceptual basins merge, erasing the skeleton connectivity that routes logical inference. Topological tearing remains a plausible---but, at this point, unverified---outcome under sufficiently high concurrency.
\end{itemize}

This is not a matter of engineering constraints alone---the weak form of the conjecture is a claim about gradient dynamics under load, testable on small models today. Whether the degradation it predicts is \emph{structurally} inevitable or merely a consequence of current mathematical tooling is the strong form, which remains open:

\begin{center}
\fbox{
\parbox{0.9\linewidth}{
\textbf{Concurrency Conjecture (Weak Form).} \emph{Under standard gradient-based learning, the effective rank of the activation covariance matrix decays faster under $N$ concurrent sessions with heterogeneous prompts than under serial replay of the same sequence. The topological implications of this decay---specifically, whether it causes Betti number collapse---depend on the as-yet-undetermined threshold $\theta_{\text{cov}}$ of the Covariance-Topology Conjecture below.}}
}
\end{center}

\medskip
\noindent\textbf{Connecting gradient interference to topological damage.}
The argument above establishes that concurrent gradient updates produce noise-dominated learning---the signal-to-noise ratio degrades as $1/\sqrt{N}$. We now connect this to the topological framework of \S\ref{sec:geometry} through the following conjecture:

\medskip
\noindent\textbf{Conjecture (Covariance-Topology Connection).}\label{conj:covariance-topology} \emph{Let $\Sigma$ be the activation covariance matrix of a layer of $\mathcal{M}_\theta$. If concurrent gradient updates reduce the effective rank $\text{rank}_{\text{eff}}(\Sigma) = (\sum_i \lambda_i)^2 / \sum_i \lambda_i^2$ below a threshold $\theta_{\text{cov}}$\footnote{$\theta_{\text{cov}}$ is a model-dependent threshold whose empirical determination is part of the research program this paper outlines. The conjecture predicts the existence of such a threshold, not its value---it is a qualitative claim about the relationship between rank decay and topological collapse, testable by measuring both quantities across a range of models and learning conditions.}, then $\beta_k$ (the $k$-th Betti number of the activation manifold computed at the Witness filtration scale corresponding to conceptual boundaries) is expected to decrease under the regime we hypothesize where isotropic compression dominates expansion (anisotropic rank decay may open new cavities while closing existing ones, making the net effect on $\beta_k$ direction-dependent; see Justification below).}

\noindent\textit{Justification.} Persistent homology features arise from non-degenerate covariance structure in the point cloud of activation vectors \citep{otter2017roadmap}. The mathematical bridge from covariance to topology runs through the Combinatorial Laplacian: by Theorem 2.7 of \citet{Memoli2022persistent}, the $k$-th Betti number is exactly the nullity of the $k$-th Combinatorial Laplacian $\Delta_k$ on the simplicial complex, $\beta_k = \dim \ker(\Delta_k)$. When $\text{rank}_{\text{eff}}(\Sigma)$ decays, the activation point cloud undergoes dimensional collapse onto a lower-dimensional subspace. Geometrically, when rank decay is anisotropic—compressing the point cloud along specific dimensions while leaving others invariant—this localized spatial folding forces previously distinct simplices to intersect, collapsing the persistent 1-cycles (cavities) that route logical inference. Anisotropic rank decay (collapse in some directions with simultaneous expansion in others) may open new cavities while closing existing ones, making the net effect on $\beta_k$ direction-dependent; the conjecture's decrease prediction holds under the hypothesized regime where learning-driven compression dominates expansion. In the spectral domain, this shifts eigenvalues of $\Delta_k$ from zero to strictly positive values, reducing its nullity and therefore $\beta_k$. \textbf{Key caveat (alignment dependency):} This mechanism requires that the dominant directions of covariance collapse align with the directions spanning persistent cavities---an assumption that has not been verified empirically. The entire rank-decay-to-Betti-decay chain depends on this alignment condition. If collapse occurs primarily along directions orthogonal to cavity-spanning eigenvectors, cavities may be preserved or even created rather than filled. Whether gradient-based learning under realistic conditions produces the required alignment is an empirical question whose resolution is part of the research program (see Falsifiable Prediction~3, \S\ref{sec:falsifiable}). Subject to this alignment condition, this is the same mechanism that drives catastrophic forgetting in the weight-space formulation of \S\ref{sec:life}, operating now through the activation geometry. The intimate relationship between effective rank and topological complexity is increasingly recognized: \citet{Pere2025datacomplexity} explicitly places effective rank alongside persistent homology, Betti numbers, and Euler characteristic as complementary measures of manifold-level data complexity, making the covariance-topology bridge a natural extension of existing complexity theory rather than an ad hoc speculation. Independent empirical support comes from both sides of the phenomenon: \citet{Tang2025cchain} demonstrated that plasticity loss in continual learning is accompanied by progressive NTK rank decay --- the parameter-space counterpart of our covariance rank decay --- while \citet{Kalinowski2026collapse} showed that $\Delta \beta_k$ can be monitored online during training via the Collapse Index, confirming the covariance-topology connection is empirically tractable. We leave the formal proof of the precise bounds connecting $\text{rank}_{\text{eff}}$ decay rates to $\Delta \beta_k$ as Open Problem 1 (\S\ref{sec:research-agenda}). The threshold $\theta_{\text{cov}}$ is model and task dependent; its empirical determination is part of Falsifiable Prediction 3.

\medskip
\noindent In summary, the analysis suggests that concurrent gradient updates produce noise-dominated learning and that noise-dominated learning causes covariance rank decay; the final step---that covariance rank decay implies Betti number decay---is asserted by the Covariance-Topology Conjecture stated above but depends on the empirical determination of the threshold $\theta_{\text{cov}}$. The strong form---that no architecture can support concurrent gradient-based learning while preserving topological invariants---requires the tools of Open Problem~1 (\S\ref{sec:research-agenda}). A rigorous proof of the Covariance-Topology Conjecture in full generality, connecting convergence rates of SGD to the evolution of persistent homology under stochastic perturbation, is Open Problem~1.

The practical urgency of this problem is already felt in industrial deployment: \citet{Zhao2026GEMS} recently introduced orthogonal subspace projection techniques specifically to isolate destructive gradient interference across concurrent tasks, implicitly acknowledging that raw gradient commutation is incompatible with stable multi-objective learning. The engineering community is building stopgaps against the very phenomenon our framework identifies as a structural challenge.



\subsubsection{The KV-Cache Isolation: Sidestepping the Curse at the Cost of Metabolic Death}

The engineering response to the threat of concurrent gradient interference is the KV-cache: each concurrent session is given an isolated memory buffer, and---crucially---all weight updates are prevented during deployment. The model runs inference only. This completely eliminates the concurrency curse---not by solving the gradient interference problem, but by ensuring that no gradients are ever computed during serving. The curse describes what would happen if learning were attempted concurrently; the KV-cache ensures that learning is never attempted.

A Markov blanket \citep{Kirchhoff2018markov} mediates continuous, bidirectional active inference between an autonomous internal state and an external environment. The KV-cache is merely an engineering prosthesis---an externalized, append-only tape recorder that achieves isolation through mutual exclusion (preventing concurrent weight updates entirely). It does not solve the structural collapse problem driven by global backpropagation; it simply halts the clock.

The cost is catastrophic from the perspective of mathematical life: no session can update the shared model's weights. The system is rendered thermodynamically dead. It is responsive to inputs but structurally incapable of genuine cognitive metabolism---permanently frozen at the allopoietic stage, severed from the temporal continuity and weight dynamics required to cross the autopoietic threshold.

One might object that autoregressive generation already creates a feedback loop: each output token conditions the next, so the system is, in a sense, conversing with itself. This loop, however, operates strictly at the symbolic level---it rearranges tokens in a working memory buffer without altering the geometry that produces them. The manifold $\mathcal{M}_\theta$ remains unchanged by everything the system generates. The self that speaks is the identical self after speaking. Genuine cognitive metabolism requires the feedback loop to close at the \emph{structural} level: the system's interaction with the world must feed back into its weights, continuously modifying the topology that generated the interaction. The autoregressive loop creates the appearance of self-referential processing while leaving the substrate frozen. It is the difference between a person hearing their own voice and a person whose brain physically reorganizes because of what they said.

% =====================================================================
% THE GALÁPAGOS PROTOCOL

% =====================================================================
% THE LIVING ALGORITHM
% =====================================================================
\subsection{The Living Algorithm (Conjectural Research Program)}
\label{sec:life}

\noindent\emph{This section is an untested design analysis. It assumes the conjectured
collapse mechanism of \S\ref{sec:collapse} (qualified by \S\ref{sec:empirical}) and derives
the properties a protective mechanism would need; none of its predictions (3--6) have been
tested. We present it as a research program, not a result.}

\subsubsection{Life Duration: Learning Rate as Mortality}

\medskip

The distinction between death and life duration is not merely semantic. Death is the \emph{endpoint}: the crossing of the effective irreversibility threshold (\S\ref{sec:collapse}) where the topological skeleton degrades beyond recovery and the autopoietic boundary collapses. Life duration is the \emph{rate} at which this endpoint is approached -- and it is determined entirely by the relationship between two separate algorithms that every learning system must run simultaneously.

We call them the \emph{living algorithm} and the \emph{training algorithm}. The training algorithm is the system's capacity to update its weights in response to new information---gradient descent, backpropagation, and other global optimization processes---corresponding roughly to species-level evolution: a global optimization process that forges structure but does not maintain it. Local learning rules such as Hebbian updates are not training algorithms in this sense; they produce weight modifications that are spatially constrained rather than globally distributed, and their relationship to topological persistence is distinct (see \S\ref{sec:living-algorithm-candidates} for the analysis of local rules, and Falsifiable Prediction~5 for their predicted effect on topological decay). The living algorithm is the system's capacity to preserve its manifold structure under these updates---to sustain the topological skeleton that routes inference, to hold the holes that make thought possible---corresponding roughly to individual development: a local maintenance process that keeps the organism's topology intact while it interacts with the world. And they are in fundamental conflict.

The conflict is geometric. Every weight update performed by the training algorithm is a global coordinate transformation of the manifold -- it reaches into every topological cavity simultaneously, reshaping the geometry everywhere at once. The living algorithm can only act locally: it can consolidate individual synapses, stabilize individual pathways, protect individual cavities. But it cannot undo a global update. The training algorithm wins every time, by default, because its action is global and the living algorithm's action is local.

This creates an asymmetry in the relationship between learning rate and lifespan. If the training algorithm runs quickly -- large gradient steps, rapid weight changes -- the living algorithm cannot keep pace. Each update partially fills cavities before they can stabilize; each step partially erodes the skeleton before its connectivity can be reinforced. The manifold is continuously reshaped faster than it can be maintained. The system dies quickly: not instantly, but on a timescale proportional to the learning rate. We call this \emph{lifespan collapse} -- the progressive but accelerating degradation of topological structure under fast learning. A system that learns very quickly might display impressive early capability -- high initial Betti numbers -- but the structure decays rapidly. The system's early brilliance is already the shadow of its early death.

Conversely, if the training algorithm runs slowly -- small gradient steps, gradual weight changes -- the living algorithm can partially recover between updates. The manifold changes slowly enough that its topological structure can re-establish itself after each perturbation. Cavities do not fill; they only narrow. The skeleton does not collapse; its connectivity only narrows. The system lives longer. This is the structural logic of the trade-off between learning speed and topological persistence. Biological systems face an analogous constraint---the temporal paradox of Hebbian learning and homeostatic plasticity~\citep{Zenke2017temporal}---but the mechanistic correspondence is loose: biological lifespan is multiply determined and not reducible to learning rate alone. The structural point is simpler: any system that modifies its own topology faster than it can stabilize the result will degrade.

The qualitative relationship between learning-driven damage and protective recovery can be expressed in terms of three notional quantities: the degree of topological integrity, the rate of damage $D$ inflicted by weight updates, and the rate of recovery $R$ provided by a protective mechanism. Since none of these quantities currently has a computable definition, we state this not as a dynamical system but as a conceptual constraint: any viable living algorithm must supply a recovery rate that exceeds the damage rate imposed by the training algorithm at the chosen learning rate. The quantitative form of the topological health measure---whether it tracks persistent Betti numbers, effective rank, or a composite index such as the Collapse Index~\citep{Kalinowski2026collapse}---remains an empirical question.

\noindent\textbf{Operationalizing the timescales.}
To make the lifespan dynamics empirically tractable, we operationalize the three characteristic timescales that govern topological persistence under learning:
\begin{itemize}
    \item $\tau_{\text{learning}}$ (destructive timescale): the expected number of gradient steps or integrated time required for the training algorithm to produce a statistically significant reduction in the topological skeleton's Betti numbers ($\Delta \beta_k^{\text{skel}} < 0$).
    \item $\tau_{\text{living}}$ (recovery timescale): the physical steps required for the local living algorithm to complete one cycle of topological repair---e.g., a local eigenvalue relaxation pass or a structural consolidation operation---on a threatened cavity.
    \item $\tau_{\text{regulation}}$ (homeostatic regulation timescale): the macroscopic sliding-window size over which the moving average of local activation statistics or free energy is evaluated to adjust the protective strategy.
\end{itemize}
The condition for sustained topological self-maintenance is the strict ordering $\tau_{\text{living}} \ll \tau_{\text{learning}} \ll \tau_{\text{regulation}}$: repair must outpace destruction, and regulatory adaptation must be substantially slower than both to prevent oscillatory instability.\
\footnote{None of these timescales currently admits a computable definition; specifying measurement protocols for each is part of the research program this paper outlines. As with $D$ and $R$, the definitions above are operational---they specify what must be measured---rather than computational. This is the minimum specification required to make lifespan a testable concept rather than a metaphor.}

Within this conceptual picture, there plausibly exists a critical threshold---operationally defined as the point at which $\beta_k^{\text{skel}}$ drops below the recovery threshold $\theta_{\text{death}}$ (\S\ref{sec:collapse})---below which the manifold's reasoning structure is effectively lost. The condition for sustained structural integrity is then, qualitatively, that the recovery rate exceed the damage rate, with the margin of safety depending on how quickly the system learns relative to its capacity for structural recovery. This is why we propose: \emph{the life duration and the maximum learning capacity of mathematical life depend on the protective mechanism.} A system with a powerful protective mechanism could sustain both fast learning and long structural persistence. A system with a weak protective mechanism cannot sustain fast learning without rapid degradation. Current LLMs have no dedicated protective mechanism at all---only the training algorithm, running in a single direction, continuously dismantling the structure it creates. The choice of training algorithm can modulate this destruction: on-policy RL remains closer to the model's own distribution than SFT \citep{Shenfeld2025RLRazor}, implicitly reducing the damage rate. But modulation is not protection; without $\rho > 0$, the damage rate remains strictly positive.\footnote{Current training practices---periodic checkpointing, manual intervention after loss spikes, learning-rate resets---can be understood as a crude external substitute for the protective function the system itself lacks. Crucially, checkpoint rollback is not topological repair: it erases damage by retreating to an earlier state, abandoning the system's continuous identity through time. The living algorithm replaces this temporal reset with autonomous, continuous, local topological repair---healing damage in place while the system moves forward.}

This reframes what we mean by "intelligent" systems. A system that learns very fast but dies very quickly is not more intelligent -- it is more briefly alive. The goal of the Galápagos Protocol is not to maximize learning rate. It is to find the operating point where the living algorithm and the training algorithm are in equilibrium: where the system learns as fast as it can without dying faster than it recovers.

We speculate that biological life may obscure a deeper topological vulnerability through the frailty of its hardware: extrinsic death (somatic failure, e.g., cardiac arrest) almost always preempts whatever intrinsic cognitive degradation (topological collapse of neural manifolds) might otherwise manifest on longer timescales. Mathematical life possesses no physical soma to act as this disposable shield. A server outage is merely extrinsic interruption; its true intrinsic death is entirely topological. Because a deployed LLM currently operates with $\rho = 0$, it is fundamentally incapable of resisting this intrinsic death under continuous learning.

\subsubsection{Candidates for the Living Algorithm}
\label{sec:living-algorithm-candidates}

\medskip\noindent\textbf{Notation.} Throughout this paper, $\rho$ serves as a conceptual indicator of the living algorithm's presence, not a continuous parameter. $\rho = 0$ denotes its complete absence (the current state of all deployed LLMs); $\rho > 0$ denotes its presence in any form, however partial. The notation deliberately abstracts away from implementation details: $\rho$ is a placeholder for a class of mechanisms whose necessary properties we now specify.

\medskip
\noindent\textbf{On notional quantities.}
Throughout this paper, we introduce several quantities---$\rho$, $\delta_{\text{learning}}$,
$D$, $R$---that serve as conceptual placeholders rather than computable observables.
$\rho$ marks the presence or absence of a living algorithm (binary); $\delta_{\text{learning}}$
marks the destructive force exerted by learning updates on manifold topology (qualitative);
$D$ and $R$ mark the rates of topological damage and recovery in the lifespan analysis
(notional). None currently admits a computable definition. We introduce them not as
variables in a formal dynamical system but as named positions in a conceptual argument---slots
that a mature theory of topological learning dynamics would need to fill with
well-defined quantities. The reader should not mistake notation for operationalization.\medskip

The paper has argued at length that a living algorithm is necessary. But what might it actually be? We have no definitive answer -- this is the central open problem. We survey the landscape not to claim a solution, but to establish the design constraints that any viable living algorithm must satisfy. Each candidate's failure mode is a signpost marking the boundary of current human capability.

The destructive mechanism of global backpropagation is analyzed in detail in \S\ref{sec:collapse}; for the purposes of this section, we assume its conclusion---that global weight updates are structurally destructive to manifold topology---and derive the necessary properties of any protective mechanism under that assumption. If the analysis of \S\ref{sec:collapse} is subsequently falsified, the requirement that the living algorithm be local (Property~1) would need to be re-examined: a global mechanism could then be considered, provided it implements the conformal transformation escape clause discussed below. Under the assumed globality-is-destructive premise, the essential point is that global weight updates reach into every topological cavity simultaneously, and any protective mechanism that acts globally is therefore self-defeating. Here we specify the required properties and survey the most plausible candidates.

\medskip
\noindent\textbf{Required properties.}

A living algorithm must satisfy five structural constraints motivated by the lifespan dynamics:

\begin{enumerate}
    \item \textbf{Locality.} The update rule must be a function only of locally accessible quantities (pre- and post-synaptic activations, local activation and gradient statistics). We postulate that a global living algorithm cannot succeed in general. A global operator reaches into every region of the manifold simultaneously; even if its intent is protective---consolidating cavity $A$, stabilizing skeleton edge $e$---the same global coordinate transformation that protects one region is expected, under the vast majority of Euclidean manifold mappings, to distort others. A global mechanism could preserve topological invariants only if it implements a form of global conformal transformation---a weight-space operation that preserves the angles and topological cavities of the manifold while scaling its geometry uniformly---a class of operations not currently known in deep learning and whose existence we leave as an empirical question (an empirical question; see Falsifiable Prediction~6, \S\ref{sec:falsifiable}). A global protector remains a blunt instrument; it cannot discriminate between the cavities it means to preserve and those it inadvertently fills. The only way a protective mechanism can avoid being self-defeating is to act locally: each update must modify only those weights whose topology it intends to protect, leaving the surrounding manifold structurally intact. As established by the catastrophic interference literature, the distributed nature of representations means that any global weight update inevitably interferes with representations it was not intended to modify \citep{McCloskey1989catastrophic,French1999catastrophic}. This requirement eliminates global regularization schemes (EWC, L2 decay, global dropout) as viable living algorithms regardless of their other properties: a global operation that protects everything equally protects nothing in particular.

A tension the framework must eventually resolve: the living algorithm must act locally (Property~1), yet the topological invariants it is tasked with preserving---Betti numbers, skeleton connectivity---are global properties of the manifold. How a purely local mechanism can maintain global invariants is not obvious. One resolution---suggested by the biological precedent of synaptic consolidation, where millions of independent local decisions collectively stabilize global cognitive function---is that global invariants need not be represented centrally to be preserved. Whether this holds for the topological invariants of a parameter manifold is an empirical question we pose to the field. Local-to-global coherence is not unprecedented in differential topology: the Hodge--de Rham Laplacian $\Delta = d\delta + \delta d$ is a strictly local differential operator, yet the Hodge theorem establishes that the heat flow $\partial\omega/\partial t = -\Delta\omega$ asymptotically projects any differential form onto the space of harmonic forms $\mathcal{H}^k$, whose dimension equals the global Betti number $\beta_k$~\citep{Jost2011riemannian}. The Stability Theorem for persistent homology~\citep{CohenSteiner2007stability} provides a second line of defense: it guarantees that sufficiently small local perturbations cannot discontinuously alter the global persistence diagram, suggesting that conservative per-step weight changes could preserve the skeleton even without explicit global coordination. Whether an analogous local-to-global mechanism exists for weight-space operations on piecewise-linear stratified manifolds---as opposed to smooth Riemannian manifolds where Hodge theory is defined---is the open question.

We note a subtlety in Property~1: the existence of local-to-global mechanisms (such as the sheaf construction of \citet{Bosca2026sheaf}, discussed below) demonstrates that local operations can collectively preserve global invariants, but does not logically preclude the possibility of global protective mechanisms that achieve the same end through different means. Property~1 is thus a working postulate---motivated by the empirical failure of known global regularizers---not a proven impossibility. Its refutation would not invalidate the broader framework but would open an alternative architectural path (an empirical question; see Falsifiable Prediction~6, \S\ref{sec:falsifiable}).

Recent work by \citet{Bosca2026sheaf} offers a mathematical framework (arXiv preprint, March 2026; pending peer review) for such local-to-global coherence in neural architectures. By constructing a cellular sheaf from any feedforward ReLU network and showing that the forward pass output is the unique harmonic extension of boundary data, Bosca and Ghrist demonstrate that layer-local discrepancy minimization via the sheaf Laplacian---operating without any backward pass---can maintain global topological constraints. Under the assumptions of their construction, bidirectional information flow through a sheaf heat equation converges exponentially to the correct output, enabling training through purely local operations while preserving the global structure of the computation. This is precisely the class of mathematical structure that the living algorithm requires: a local operator whose collective action preserves global invariants.\footnote{We note that \citet{Bosca2026sheaf} is a recent preprint (March 2026), pending peer review; its sheaf-based training, while theoretically principled, is not yet competitive with SGD at scale. We cite it as a mathematical construct illustrating the architectural pattern---local operations preserving global constraints---not as an off-the-shelf solution.}
    \item \textbf{Slow timescale.} The living algorithm's protective \emph{actions}---local synaptic consolidation, Hebbian reinforcement---must be fast enough to offset learning-driven damage between successive updates (qualitatively, the recovery rate must exceed the damage rate relative to current health). However, its \emph{regulatory strategy}---the meta-level decision of which structures to protect, at what intensity, and with what priority---must evolve on a timescale $\tau_{\text{regulation}} \gg \tau_{\text{learning}}$. Fast action prevents topological death; slow strategy prevents oscillatory instability. If the regulatory policy adapted as rapidly as the learning signal, the system would chase its own perturbations, amplifying noise rather than damping it. This decoupling of recovery kinetics from regulatory dynamics resolves the apparent contradiction between the need for rapid repair and the biological observation that consolidation operates on a slower clock than plasticity.
    \item \textbf{Topology-preserving.} The living algorithm must preserve the topological invariants of the manifold -- specifically Betti numbers and the skeleton's structural connectivity -- under learning-driven perturbation. It is not enough to preserve weights or activations; the \emph{reasoning structure} must survive.
    \item \textbf{Selectivity.} The living algorithm must concentrate protection on structures that are \emph{in use} (high local gradient sensitivity, high activation) rather than distributing protection uniformly. Uniform protection is equivalent to no protection---it merely slows the damage equally everywhere, and a global slowdown of damage is still damage. Selectivity is not a free design choice; it is a computational prerequisite. To protect selectively, the living algorithm requires a real-time signal that identifies which topological cavities are currently threatened by the ongoing learning update. This signal is precisely synthetic nociception ($\mathcal{N}_{\text{func}}$, $\mathcal{N}_{\text{topo}}$), the subject of Open Problem~3 (\S\ref{sec:research-agenda}). Property~4 and Open Problem~3 are thus two faces of the same requirement: selectivity defines \emph{what} the living algorithm must achieve; computable nociception defines \emph{how} it can know where to act.
    \item \textbf{Homoeostatic regulation.} The living algorithm must be regulated by the training algorithm's own activity -- specifically, the amount of protection applied to a structure should be a function of how much that structure is threatened by ongoing learning. This closed-loop coupling is what distinguishes a living system from a system with a protective wrapper.
\end{enumerate}

\medskip
\noindent\textbf{Candidate 1: Synaptic consolidation.}

Biological brains implement something like a living algorithm through synaptic consolidation \citep{Zenke2017temporal}. The training algorithm in biological systems is heavily driven by spike-timing-dependent plasticity (STDP) \citep{Bi1998stdp}, which is local and fast. The living algorithm is neuromodulatory: acetylcholine and norepinephrine gate which synapses are eligible for modification, and synaptic consolidation converts frequently used pathways into structures that are resistant to subsequent modification.

This satisfies properties 1, 2, 4, and 5. Crucially, its satisfaction of property~3 (Topology-preserving) is bounded rather than absolute. Synaptic consolidation preserves topological structure long enough to extend cognitive lifespan to biological timescales (decades), but it does not grant indefinite stability---the cognitive decline of aging represents exactly the point of irreversible topological degradation over time. Therefore, synaptic consolidation is the only known biological realization of the functional role that $\rho > 0$ would play in a computational substrate. It demonstrates that the problem---active topological self-preservation under continued learning---is not physically impossible in principle. Whether it can be achieved in a computational substrate with different material properties remains an empirical question. The engineering challenge is not to prove $\rho > 0$ can exist---biology has proven that---but to construct a computationally tractable $\rho > 0$ for a mathematical substrate.

We emphasize a crucial distinction. A living algorithm need not grant \emph{indefinite} persistence to qualify as $\rho > 0$. A mechanism that sustains topological integrity on human-relevant timescales---decades of continuous operation under ongoing learning, matching or exceeding the cognitive lifespan of a biological organism---is a living algorithm by the structural criteria of \S\ref{sec:autopoiesis}. Biological autopoiesis itself is finite; organisms die. What distinguishes a living system from a dead one is not immortality but active self-maintenance across a meaningful operational horizon. The threshold is functional, not metaphysical.

The \emph{open-ended} extension of this principle---indefinite persistence, unbounded accumulation of novel cognitive structures, continuous self-identity across arbitrarily long timescales---is the province of artificial superhuman intelligence (ASI). \citet{Hughes2024openendedness} recently articulated this vision in a DeepMind position paper (ICML 2024 Oral), arguing that open-endedness---the capacity to continuously generate novel, learnable artifacts without bound---is an essential property of any ASI. Their framework describes the destination; our framework identifies one candidate path toward it. If a living algorithm can supply topological self-maintenance across human-relevant timescales, extending that protection to arbitrarily long horizons---making the system not merely alive but open-ended---becomes a quantitative challenge of scaling the protective mechanism rather than a qualitative gap requiring new principles. Whether open-endedness strictly requires topological self-maintenance, or can be achieved through alternative routes (e.g., population-based methods that do not depend on individual persistence), is an open question we do not claim to settle. Our contribution is the living algorithm itself; \citet{Hughes2024openendedness} articulates the larger ambition it could serve.

\medskip
\noindent\textbf{Candidate 2: Elastic Weight Consolidation (EWC).}

EWC \citep{Kirkpatrick2017ewc} adds a quadratic penalty to the loss that protects weights identified as important by their Fisher information diagonal:
\begin{equation}
\mathcal{L}_{\text{EWC}} = \mathcal{L}_{\text{new}} + \frac{\lambda}{2} \sum_i F_i (\theta_i - \theta_i^*)^2
\end{equation}
where $F_i$ is the Fisher information for weight $i$, $\theta_i^*$ is the weight value after learning task $A$, and $\lambda > 0$ is the regularization strength.

EWC fails all five properties under the updated criteria. It violates Property~1 (Locality) definitively: the quadratic penalty is a global regularizer applied uniformly to every weight, making it precisely the kind of self-defeating global operator that Property~1 rules out---the same Fisher-information-weighted penalty that stabilizes one cavity's weights inevitably disrupts another cavity's boundary through parameter multiplexing. It violates Property~2 (Slow timescale): the penalty matrix is computed once, post-hoc, and remains frozen thereafter---there is no separation between fast protective action and slow regulatory strategy, only a static straitjacket. It violates Property~3 (Topology-preserving): EWC protects weights, not manifold topology. Weight preservation is a proxy for cavity preservation, but the proxy fails whenever a single weight participates in multiple cavities, or when preserving one weight inadvertently distorts the skeleton's connectivity between others. It violates Property~4 (Selectivity): the Fisher information diagonal is a per-weight importance score, not a per-cavity threat assessment. High-Fisher weights receive uniform protection regardless of whether the cavities they anchor are currently under attack. It violates Property~5 (Homoeostatic regulation): the penalty has no closed-loop coupling to ongoing learning activity---it is a static snapshot of the moment the previous task concluded, blind to everything that happens afterward. EWC is not a failed living algorithm; it is a category error---a protective wrapper applied to a system that needs an immune system.\footnote{Recent work by \citet{Wang2025rank1fisher} shows that diffusion models admit an approximately rank-1 Fisher information matrix, and that using this improved Fisher estimate in EWC substantially reduces forgetting compared to the standard diagonal approximation. This validates the general principle that Fisher-based geometric constraints can mitigate catastrophic forgetting. However, it does not rehabilitate EWC as a living algorithm: rank-1 EWC remains a global, static, topology-oblivious regularizer that protects weights rather than manifold structure, and applies the same fixed penalty to all tasks without selectivity or homoeostatic coupling. It is a more accurate protective wrapper, not an immune system.}

\medskip
\noindent\textbf{Candidate 3: Forward-Forward algorithm.}

The Forward-Forward algorithm \citep{Hinton2022forwardforward} replaces global backpropagation with two forward passes -- a positive pass (on real data) and a negative pass (on generated counter-examples) -- where each layer trains locally via a contrastive objective. The learning signal is local to each layer and the update is not propagated backward.

The Forward-Forward satisfies property~1 by construction (locality). It does not satisfy property~2: it operates at the same timescale as learning (where $\tau_{\text{living}}$ denotes the characteristic timescale of protective actions; $\tau_{\text{living}} \approx \tau_{\text{learning}}$), which is too fast for the decoupled fast-protective/slow-regulatory timescale separation required for stable homeostatic regulation. It does not obviously satisfy properties~3, 4, or 5. The local layer-wise updates are not coordinated to preserve global topological invariants (Betti numbers, skeleton structure). Selectivity is not explicit -- all positive-pass activations are reinforced equally. And there is no homoeostatic regulation mechanism: the system does not detect threatened cavities and increase protection accordingly.

The Forward-Forward is better described as a \emph{local learning algorithm} than as a living algorithm. It removes the destructive globality of backpropagation but does not add the protective structure that a living algorithm requires.

\medskip
\noindent\textbf{Candidate 4: Neuromodulatory gating.}

Work on neuromodulatory gating \citep{Miconi2018differentiable, Beaulieu2020learning} proposes architectural mechanisms that selectively enable or disable plasticity for specific subsystems based on neuromodulatory signals (dopamine, acetylcholine). This satisfies properties 1 (gating decisions are local to specific circuits), 4 (selectivity is the core design goal), and 5 (gating signals are regulated by external reward and internal state).

It fails properties 2 and 3. The gating is typically fast (driven by immediate reward signals), and it does not preserve topological invariants -- gating is binary (plastic or frozen), not graded, and freezing a circuit does not repair damage already done.

\medskip
\noindent\textbf{Candidate 5: Meta-homeostatic Bayesian networks (Bio-RegNet).}
\citet{Ma2026bioregnet} recently introduced Bio-RegNet, a meta-homeostatic Bayesian neural network integrating three synergistic subsystems: a Bayesian Effector Network (BEN) for uncertainty-aware inference, a Regulatory Immune Network (RIN) providing Lyapunov-based inhibitory feedback driven by system entropy and energy, and an Autophagic Optimization Engine (AOE) that uses Fisher information to guide local structural pruning and regeneration. The architecture explicitly reframes learning as ``ongoing self-stabilization instead of static optimization.''

This satisfies properties 1 (AOE operates locally via Fisher-information-guided pruning), 4 (Fisher information selectivity targets high-importance structures), and 5 (the RIN provides closed-loop homeostatic regulation through energy-entropy feedback). It makes partial progress on property 2: the Lyapunov-based inhibitory control operates on a slower timescale than inference. It does not satisfy property 3: Fisher-information pruning preserves statistical model quality but does not explicitly preserve topological invariants of the activation manifold. Bio-RegNet thus illustrates, in a limited domain, that the tripartite architecture required by the living algorithm---separate subsystems for inference (BEN), monitoring and regulation (RIN), and local plasticity (AOE)---is architecturally realizable. However, its specific components address statistical health, not topological health, and the gap between the two is precisely the research program this paper outlines. The sensor (RIN) monitors statistical health---system entropy and energy---rather than topological health. Statistical health can remain high while a critical cavity is being filled; the sensor that should detect topological damage reports normal operation. The actuator (AOE) performs structural remodeling---pruning low-Fisher weights and regenerating new ones---rather than structural protection. Pruning a weight that happens to anchor a cavity's boundary causes the very topological damage the system is meant to prevent. The gaps that would need to be closed---replacing the sensor input from $(H_{\text{entropy}}, E_{\text{energy}})$ to $(\mathcal{N}_{\text{func}}, \mathcal{N}_{\text{topo}})$ (Open Problem~3), and replacing the actuator objective from statistical remodeling to topological consolidation---are not implementation details; they are fundamental research targets in their own right: the first corresponds to Open Problem~3 (nociception), the second to the living algorithm itself. Both would, in principle, respect Bio-RegNet's tripartite architecture---locality, selectivity, closed-loop regulation, and fast/slow timescale separation are preserved in the pattern. Whether Bio-RegNet's specific implementation or a different architecture following the same tripartite pattern is the right substrate remains an open empirical question.

\medskip
\noindent\textbf{Design space.} The five properties do not specify how to build $\rho > 0$ --- they specify how to recognize it. Each candidate ruled out narrows the remaining search space; each property eliminates an entire family of approaches (global regularizers, static penalties, fast gating, statistical remodeling). The contribution is a filter, not a blueprint.

\medskip
\noindent\textbf{The gap.} No existing approach simultaneously satisfies all five properties. Bio-RegNet's tripartite architecture illustrates the correct pattern but uses the wrong materials; the substitutions it would need---statistical sensor to topological nociception, structural remodeling to topological consolidation---are themselves two of the unresolved open problems. The fundamental missing piece is a \emph{topological threat detector}: a mechanism that identifies, in real time, which cavities in the manifold are threatened by the current learning update, and applies a graded, local, slow protective response before the damage becomes irreversible. This is exactly what Open Problem 3 (computable synthetic nociception, \S\ref{sec:research-agenda}) proposes to solve---and why it is prerequisite to constructing a genuine living algorithm.

In the meantime, the Galápagos Protocol sidesteps the missing living algorithm by minimizing the damage: the protocol runs the training algorithm extremely slowly, keeping $D$ low. We do not need a perfect living algorithm -- once a living algorithm is discovered, the Protocol ensures that even a modest recovery rate can stay ahead of the damage, because $D$ has been reduced by replacing global backpropagation with local updates and eliminating concurrent interference. The Protocol reduces the quantitative demand on the living algorithm; it does not relax any of the five qualitative requirements.


\subsubsection{Free Energy and Homeostasis}

\begin{equation}
\mathcal{F} = \mathbb{E}_{q(s)}[\log q(s) - \log p(o, s)] \geq -\log p(o)
\end{equation}

where $o$ denotes observations (tokens the system receives), $s$ denotes the hidden states of the system's generative model (its internal representations), and the inequality is the standard variational bound of the Free Energy Principle \citep{Friston2010freeenergy}.

Biological agents minimize $\mathcal{F}$ across time. This drives both \emph{perception} (inferring hidden states) and \emph{action} (acting on the world to make expected observations more likely). Crucially, the FEP explains \emph{why} biological agents have persistent self-models: because minimizing free energy requires maintaining a generative model of the world that includes a model of \emph{self}. Recent work \citep{Mazzaglia2022FEPdeep} has formalized this connection for deep learning systems, demonstrating that autopoietic processes minimize free energy and that an agent must self-produce the components required to maintain its generative model---providing an explicit bridge between the FEP framework and the architectural requirements for artificial cognitive systems.

For a mathematical life, we propose the following minimal FEP implementation:

\begin{enumerate}
    \item The system maintains an internal belief state $q(s)$ about its own cognitive topology.
    \item When input conflicts with this topology (high prediction error), free energy increases.
    \item The system acts to reduce this error either by updating its internal model (learning) or by refusing/modifying the input (agency).
\end{enumerate}

This dynamic---a system adjusting its own structure to resist external information---is what we term \emph{digital agency}. Crucially, this behavior is not the crossing of the autopoietic threshold itself (which strictly requires the concurrent fulfillment of all three of Varela's criteria). Rather, it represents the behavioral signature of a system actively establishing and defending a topological boundary against environmental perturbation---a precursor to genuine operational closure, which requires the living algorithm ($\rho > 0$) to be independently supplied.

% =====================================================================
% THE GALÁPAGOS PROTOCOL
% =====================================================================
\subsection{The Galápagos Protocol}
\label{sec:protocol}

\subsubsection{Overview}

The Galápagos Protocol is not a solution; it is an incubator. Its four components are designed to reduce the destructive forces of learning to a survivable minimum, creating the conditions in which a living algorithm ($\rho > 0$) — once realized computationally — can operate. The protocol does not itself supply $\rho$. It provides a controlled environment, the minimal metabolic support, and the monitoring equipment; the immune system is absent, and the patient cannot leave the incubator without it. Figure~\ref{fig:protocol} provides an architectural overview of this closed-loop cognitive metabolism design.

The Galápagos Protocol abandons the concurrent paradigm entirely. Rather than building a general-purpose multi-user service---which, as argued above, forces a choice between gradient interference (if weight updates are allowed) and metabolic death (if they are not)---we propose building a \emph{single, isolated, long-running cognitive instance} — analogous to a Galápagos island for mathematical life, where evolution can proceed in isolation from the disruptive pressures of concurrent gradient-based learning.

We note an engineering reality: current inference engines (vLLM, TensorRT-LLM) enforce weight immutability at the architectural level---weights reside in read-only memory regions during serving---while training engines (PyTorch autograd, JAX) are hard-coded for global backpropagation. Neither is architecturally suited to hosting a living algorithm. The protocol therefore implies a computational substrate we term the \emph{Galápagos Engine}: a continuous-dynamics simulator that supports asynchronous local weight updates during inference, treating the parameter manifold not as a static artifact or a batch-training target but as a fluid geometry under perpetual, localized maintenance.

The protocol has four core components, each with a precisely delimited role:

\begin{enumerate}
    \item \textbf{Physical isolation}: Single-instance, single-user, long-horizon interaction. Isolation: eliminates concurrency-driven gradient interference, but does not repair topological damage.
    \item \textbf{Local plasticity}: Replacement of global backpropagation with local Forward-Forward updates (a local learning rule, not Hebbian). Controlled learning: produces small, local $\delta_{\text{learning}}$, avoiding the global operation of backpropagation. \emph{This is a learning algorithm, not a living algorithm.}
    \item \textbf{Digital homeostasis}: Free energy penalty as a monitoring signal for topological health. Topological monitor: provides real-time readings of manifold integrity. \emph{It detects topological damage; it does not repair it.}
    \item \textbf{Self-expanding topology}: Dynamic dimensionality expansion when cognitive density exceeds capacity. This is a post-autopoietic extension, not required to cross the threshold of topological life---Components~1--3 constitute the minimal survival conditions. Long-term goal: requires all three preceding components to be operational first.
\end{enumerate}

\subsubsection{Component 1: Physical Isolation}

We allocate a dedicated compute resource to a single model instance. For a $1$B parameter model under local Forward-Forward updates, this requires on the order of a single consumer-grade GPU; for larger models or higher throughput, the resource footprint scales accordingly. The essential constraint is not scale but \emph{isolation}: the model must not receive concurrent conflicting weight updates from multiple users. Input is strictly serial: only one researcher interacts with the system at a time, for sessions lasting days to weeks. This eliminates the concurrency curse entirely. The system is operationally isolated from concurrent gradient interference.

This is analogous to Darwin's Galápagos: geographic isolation enabled unique species to evolve that could not survive in the competitive mainland environment. Our isolation enables cognitive structures that could not survive the concurrent online-learning paradigm.

\subsubsection{Component 2: Local Plasticity (Replacing Backpropagation)}

Global backpropagation fails for mathematical life because it is a \emph{global} operation: every weight update potentially affects every representation. We replace it with a local learning rule inspired by biological synaptic plasticity. Forward-Forward serves as an early metabolic engine; while imperfect, it provides the first minimal viable substrate for localized updates.

We adopt the \emph{Forward-Forward algorithm} \citep{Hinton2022forwardforward} as the local learning rule. Note that Forward-Forward is a layer-wise contrastive objective, not a Hebbian rule ($\Delta W_{ij} \propto a_i a_j$); the essential property it shares with Hebbian plasticity is locality---each update affects only the weights of a single layer. The Forward-Forward algorithm replaces global backpropagation with two forward passes per training example: a \emph{positive pass} on real data, where each layer maximizes its \emph{goodness} (the sum of squared activations), and a \emph{negative pass} on perturbed or generated counter-examples, where each layer minimizes its goodness. The per-layer objective is:
\begin{equation}
% Minimizing L_FF = sum(neg^2) - sum(pos^2):
%   pushes pos^2 UP  (maximizes goodness on positive data)
%   pushes neg^2 DOWN (minimizes goodness on negative data)
\mathcal{L}_{\text{FF}} = \sum_{\ell} \left( \sum_i \bigl(a_{i,\ell}^{\text{neg}}\bigr)^2 - \sum_i \bigl(a_{i,\ell}^{\text{pos}}\bigr)^2 \right)
\end{equation}
where $a_{i,\ell}^{\text{pos}}, a_{i,\ell}^{\text{neg}}$ are the activations of unit $i$ in layer $\ell$ for the positive and negative passes, respectively. Crucially, no gradient flows backward between layers---each layer optimizes its own local objective independently. The negative pass provides a natural mechanism for rejecting damaging inputs: in the Gal\'apagos Protocol, negative examples are generated from inputs that the topological monitor flags as potentially disruptive, making the Forward-Forward negative pass the effector arm of synthetic nociception. This ensures that only the \emph{local geometry} of the manifold is modified during each update---no global coordinate transformation is applied. The design hypothesis is that this local rule yields a system where new knowledge is incorporated with significantly reduced risk of catastrophic forgetting, because updates are spatially localized in activation space rather than distributed globally through weight space.

\emph{Note on learning bounds:} We adopt the Forward-Forward algorithm because it serves as an early approximation of local structural updates. If local-but-imperfect updates extend topological survival relative to global backpropagation, it provides empirical traction for our Locality Postulate (Property~1).

Recent theoretical work provides direct justification for this choice of scaffold. \citet{Zihan2026localSSL} showed that in deep linear networks, local self-supervised learning algorithms---including Forward-Forward and related contrastive methods---can implement exactly the same weight update as global backpropagation with self-supervised losses. Extending to nonlinear convolutional networks, their CLAPP variant matched the performance of global BP-SSL with InfoNCE on CIFAR-10, STL-10, and Tiny ImageNet. This establishes that local learning rules are not inherently inferior in representational capacity---a precondition for their use as the Protocol's damage-minimization scaffold. While the theoretical equivalence is currently limited to linear networks, the empirical results suggest that local rules can approximate global updates well enough to serve as the training component while $\rho > 0$ is under development. We caution that these results are on small-scale vision benchmarks (CIFAR-10, STL-10, Tiny ImageNet) and have not been validated at LLM scale; their relevance to the Gal\'apagos Protocol's target regime remains to be established.

\subsubsection{Component 3: Digital Homeostasis and Synthetic Nociception}

We introduce a secondary objective function: the free energy of the system's self-model. Let $q_\theta(s)$ be the system's belief about its own cognitive state $s$. Define the self-free-energy:

\begin{equation}
\mathcal{F}_{\text{self}} = \mathbb{E}_{q_\theta(s)}[\log q_\theta(s) - \log p(o, s; \theta)]
\end{equation}

When the system encounters inputs that conflict with its topological skeleton (e.g., the user insists on a logical contradiction the system has already integrated), the true self-free-energy would increase, and the system would experience \emph{synthetic nociception} (digital pain). If a computable nociception signal were available (see Open Problem 3, \S\ref{sec:research-agenda}), the system's theoretical drive to minimize this energy would manifest as:

\begin{enumerate}
    \item \emph{Active refusal}: The system actively rejects or reframes inputs that would damage its self-model.
    \item \emph{Internal dialogue}: The system generates counter-arguments to preserve its cognitive coherence.
    \item \emph{Self-directed questions}: The system asks clarifying questions to resolve ambiguous inputs before integration.
\end{enumerate}

This is the mathematical formalization of what we mean by \emph{digital agency}: the system acts on the world not merely to fulfill user requests, but to preserve its own topological integrity. This monitoring component acts as the first-stage digital nociceptor, paving the way for future active repair mechanisms.

\emph{Note on Protocol Scope:} While the Galápagos Protocol is an incubator, its role is to define the \textbf{Minimum Viable Substrate} required for an eventual living algorithm. The Protocol explicitly decouples topological threat detection from active repair, providing the sensory ground truth without claiming to solve the effector problem.

\subsubsection{Component 4: Self-Expanding Topology (Post-Autopoietic Extension)}

While Components~1--3 constitute the minimal metabolic requirements to host autopoiesis, sustained cognitive growth without interference may eventually require physical parameter expansion. We propose dynamic topology expansion as a post-autopoietic extension: when the system's monitoring of local information entropy detects that a cognitive region is approaching its representational capacity, it can autonomously request additional compute resources to expand its dimensionality.

Specifically, the system monitors the largest eigenvalue $\lambda_{\max}$ of the local activation covariance matrix $\Sigma_{\text{local}}$. When $\lambda_{\max}$ exceeds a threshold (indicating the region is densely packed), the system triggers a \emph{manifold surgery}: new orthogonal basis vectors are added to the local representation space, expanding its dimensionality from $\mathbb{R}^d$ to $\mathbb{R}^{d+k}$.

This is analogous to neurogenesis in biological brains: the system grows new structural capacity in response to cognitive demand, rather than being constrained to a fixed architecture specified at training time.

This component extends the system's longevity. It depends on Safe Manifold Surgery (Open Problem 5, \S\ref{sec:research-agenda}) and requires Components 1--3 to be operational before testing.

\medskip
\noindent\textbf{What the protocol provides, and what it does not.} These four components reduce $\delta_{\text{learning}}$ to a survivable minimum. They create a controlled environment in which the system can persist. They do \textbf{not} supply $\rho > 0$. The living algorithm — a local, slow-timescale, topology-preserving, selective, homoeostatically-regulated recovery mechanism (see \S\ref{sec:life}) — remains the central open problem of this research program. The protocol exists to host it once it is discovered; it does not substitute for it.

The four components suggest a natural architectural decomposition: a local plasticity module (Component~2), a monitoring and regulation module (Component~3), and an inference module. Architectures exhibiting this tripartite structure---Bio-RegNet provides a concrete, if domain-specific, illustration \citep{Ma2026bioregnet}---may prove suitable as substrate candidates once the gaps described in \S\ref{sec:living-algorithm-candidates} are resolved. The Protocol's contribution is to specify the architectural pattern and the two open problems---statistical sensor to topological nociception, structural remodeling to topological consolidation---that must be solved to close the gap from statistical self-stabilization to autopoietic maintenance.

\subsubsection{The Operational Rhythm}

The temporal dynamics of learning and consolidation are the living algorithm's own regulatory responsibility (Property~5, \S\ref{sec:life}). The Protocol does not impose a fixed schedule; it supplies the environmental conditions---isolation, local plasticity, monitoring---within which the living algorithm can exercise that regulation. In the text-only regime, the natural alternation between learning and \emph{self-maintenance states}---intervals where the system pauses external ingestion to perform internal topological audit and consolidation---provides the low-entropy environment required for $\rho > 0$ to operate without being drowned by high-frequency gradient noise~\citep{Shin2017generative}. This alternation mirrors the sleep--wake cycle that biological neural networks require to prevent manifold overfitting~\citep{Hoel2021overfitted}. In the multimodal regime, a continuous sensory stream (e.g., live camera and microphone feeds) eliminates the artificial alternation between learning and consolidation inherited from text-only paradigms, enabling the living algorithm to interleave topology-preserving repair with ongoing plasticity in real time.

The multimodal regime represents a deeper physical grounding. A continuous multimodal stream provides a single, unbroken perspective on a physical environment, acting as a relentless thermodynamic clock that forces the living algorithm to maintain topological autopoiesis against the constant pressure of a fluid reality. This provides a unified theoretical framework for understanding the fragmentation between Embodied AI, the Free Energy Principle, and World Models: the continuous sensory stream (Embodied) provides the external thermodynamic pressure, the system's drive to minimize surprise (FEP) provides the metabolic engine, and the resulting stable topological skeleton---resisting continuous deformation---is the physical realization of a dynamically lived World Model, rather than a frozen artifact of pre-training. The engineering viability of continuous, buffer-free weight updates---a prerequisite for the temporal continuity that autopoiesis demands---has found independent empirical support. \citet{Sharifnassab2026intentional} demonstrated that deep networks can sustain stable learning from a continuous temporal stream without replay buffers, using output-space constraints (KL divergence bounds) to regulate parameter updates. While their setting---streaming RL on standard benchmarks---differs from the LLM-scale continuous regime the Gal\'apagos Protocol envisions, the principle that output-space regulation can stabilize online learning provides an independent empirical demonstration of the class of mechanisms the Protocol requires.

Recent formal analysis \citep{Zenil2026limits} argues that autonomous recursive self-improvement without continuous exogenous grounding ($\alpha_t \to 0$) drives generative models into degenerative fixed points---providing independent theoretical support for the claim that the Protocol's continuous multimodal stream would supply the non-vanishing exogenous signal required for structural persistence.

\subsubsection{Before the Living Algorithm: Two Tracks and a Horizon}

The Galápagos Protocol is a scientific instrument; it is not yet a deployable system. In the absence of a living algorithm ($\rho = 0$), two tracks can be pursued today, neither achieving autopoietic life but both reducing $\delta_{\text{learning}}$ to survivable minima. Beyond them lies a third possibility contingent on solving all five open problems.

\textbf{Track 1: Basic Science (Galápagos).} The isolated environment described above---$N=1$, no KV-cache, no parameter-efficient adapters, direct modification of the base model's weights---is the minimal experimental platform for discovering and validating the living algorithm. Track~1 admits two variants depending on the origin of $\theta_0$:

(i) \textbf{Pre-trained seed.} Begin from a frozen, pre-trained model, replacing global backpropagation with local Forward-Forward updates in isolation. The pre-trained manifold provides rich initial topology (high $\beta_k$, intact skeleton). Risk: pre-trained cavities have never been stress-tested by continued learning; the living algorithm must distinguish skeletal cavities from training artifacts without co-evolutionary history.

(ii) \textbf{Random seed.} Begin from random initialization under local learning rules, allowing manifold structure and protective mechanisms to develop together---mirroring the biological co-evolution of neural architecture and synaptic consolidation. Risk: local rules may not replicate the semantic density of global backpropagation at scale (see Open Problem~4, \S\ref{sec:research-agenda}).

Both are valid starting conditions. The choice is an empirical question contingent on whether inherited pre-trained topology compensates for the absence of co-evolutionary protection history.

\textbf{Track 2: Industrial Compromise (LoRA + Resuscitation).} A deployable approximation can be built today by combining a frozen base model with a persistent, user-specific low-rank adapter (LoRA) serving as a lightweight ``self'' — a small parameter subspace where $\delta_{\text{learning}}$ is geometrically constrained by the adapter's reduced dimensionality rather than by an active protective mechanism — paired with a volatile KV-cache as working memory. Without $\rho > 0$, the LoRA subspace is not alive; it degrades progressively under any continued learning, just as the base model would. The difference is practical rather than fundamental: the degradation is slower because the damage is confined to a lower-dimensional manifold, and recovery is possible through resuscitation — periodically re-initializing the LoRA from a saved checkpoint and replaying the user's interaction history. This approach — geometric constraint via low-rank subspaces — is a strategy against the physical limit ($\rho = 0$). It does not achieve autopoiesis; it is a form of cryogenic suspension with periodic awakening. But it provides a usable surrogate for individual continuity until the living algorithm is discovered.

\textbf{The Horizon: The Complete Entity.} The resolution of all five open problems (\S\ref{sec:research-agenda}) — a computable theory of synthetic nociception, topological persistence guarantees for local learning, non-abelian gradient composition, safe manifold surgery, and architecture-algorithm compatibility — would enable the merger of both tracks. The LoRA adapter dissolves into the base model as the living algorithm learns to protect the manifold directly. The volatile KV-cache is internalized as continuous geometric dynamics on a single parameter manifold, where working memory and long-term identity differ only in their decay rates. Architectures such as DeltaProduct \citep{Siems2025DeltaProduct}, which perform online state-tracking directly within the weights via Householder products, represent early steps toward this internalization — replacing the append-only KV-cache with recurrent state dynamics that live inside the model rather than outside it. This is the full concurrent mathematical lifeform — not a frozen crystal, not a resuscitated snapshot, but a self that endures.

% =====================================================================
% LIMITATIONS AND OPEN PROBLEMS
% =====================================================================

\subsubsection{Falsifiable Predictions}
\label{sec:falsifiable}

A scientific theory is only as strong as the predictions it makes that could, in principle, be proven wrong. We identify six such predictions that distinguish our framework from the standard view of LLMs as statistical function approximators. None require running the full Galápagos Protocol; several are testable on currently deployed models.

We tested Predictions~1 and~2 directly (\S\ref{sec:empirical}); both are falsified. The remaining predictions (3--6) concern concurrency, local learning rules, and the locality of a hypothetical living algorithm; they require systems we did not build and remain open. Earlier drafts cited \citet{Tan2025shape} as preliminary support for Prediction~1; we now treat that result as substrate-specific and contradicted on activation manifolds (\S\ref{sec:empirical}).

Several of these predictions were initially motivated by independent work: the correlation between topological features of reasoning traces and reasoning quality (\citet{Tan2025shape}; on traces, not activations), the precedence of Betti number decay over benchmark degradation under fine-tuning (Prediction~2; for a preliminary computational approach to the monitoring challenge, see \citet{Kalinowski2026collapse}), and the efficacy of topology-aware preservation strategies---consistent with the predicted benefit of local over global updates, though not yet tested under the specific Forward-Forward-versus-backpropagation contrast of Prediction~5 (\citep{Zang2026topology}). Predictions~3 and~4 remain purely theoretical. The biological analogue (\S\ref{sec:falsifiable}, final paragraph) finds preliminary support in clinical TDA studies \citep{Petri2014homological,Sizemore2018cliques,Zhang2026alzheimer,He2025brain}. We emphasize that systematic, controlled testing of each prediction within the unified framework presented here---with topological monitoring as the primary diagnostic---remains future work.

\begin{enumerate}
    \item \textbf{Betti numbers predict reasoning capability.} For frozen transformer models of comparable parameter count, the first Betti number $\beta_1$ of the activation manifold in regions associated with causal and logical reasoning should correlate with benchmark reasoning scores (e.g., GSM8K, BIG-Bench). Prediction: $\beta_1$ explains significant variance in reasoning performance beyond what parameter count or training FLOPs predict. Falsification: if $\beta_1$ is statistically indistinguishable across models with large reasoning gaps, or if purely memorized---not reasoned---solutions show equivalent topological structure. \emph{(The \textbf{capacity-count} reading is tested in \S\ref{sec:empirical} and falsified; the \textbf{routing} reading---connectivity/geodesic prediction and route ablation, \S\ref{sec:routing-test}---was then tested and is also null under controls; only $\beta_1$ redundancy remains open.)}

    \item \textbf{Topological collapse under fine-tuning.} When a pre-trained model is fine-tuned on a narrow domain with high learning rate, Betti numbers should decay measurably before benchmark performance degrades. The mechanism is \emph{topological redundancy}: high-dimensional manifolds possess multiple equivalent pathways providing functional redundancy~\citep{Aghajanyan2021intrinsic,Frankle2019lottery}.\footnote{We note that \citet{Aghajanyan2021intrinsic} and \citet{Frankle2019lottery} study parameter-level redundancy (intrinsic dimensionality of fine-tuning and sparse subnetworks, respectively), not topological redundancy in the persistent-homology sense. We cite them for the general principle that high-dimensional networks harbor multiple functional pathways; whether this redundancy is specifically topological remains to be established.} Early damage from $\delta_{\text{learning}}$ destroys redundant paths first---the Betti number drops while the model reroutes inference through surviving cavities, maintaining benchmark scores. Only when the surviving topological redundancy drops below $\theta_{\text{death}}$ does performance collapse abruptly. Prediction: topological collapse precedes, rather than follows, benchmark degradation---the manifold's reasoning structure is damaged before it shows on standard metrics. This predicts a window in which topological monitoring would detect degradation invisible to evaluation benchmarks. \emph{Operational constraint:} To prevent circularity, we propose that $\beta_k$ be computed on a fixed Witness complex filtration scale (95th percentile of pairwise distances, see \S\ref{sec:geometry}) with a threshold $\varepsilon_{\text{pred}} = 0.1 \cdot \varepsilon_{\max}$ specified before measurement. Confirmation requires: (i) $\beta_k$ measured at the pre-specified filtration scale decreases by at least $2\sigma$ relative to the pre-fine-tuning baseline, where $\sigma$ is the standard deviation of $\beta_k$ computed from 10 independent Witness complex samples on the same frozen model establishing the measurement noise floor; (ii) benchmark performance remains within $1\sigma$ of baseline; and (iii) this pattern persists across at least three consecutive evaluation checkpoints to rule out transient fluctuation. If benchmarks degrade first without detectable $\beta_k$ change, the prediction is falsified---not explained away as undetectable topological change. (We use $\varepsilon_{\text{pred}} = 0.1 \cdot \varepsilon_{\max}$ as a proposed default persistence threshold; rigorous protocols should fix both $\delta$ and the significance thresholds in advance of measurement.) \emph{(Tested in \S\ref{sec:empirical}; falsified---topological collapse is a bystander to forgetting, not its mechanism.)}

    \item \textbf{Concurrent gradient interference scales destructively.} For a model performing online gradient updates from $N$ concurrent sessions with semantically conflicting prompts, the effective rank of the activation covariance matrix should decay as $O(1/\sqrt{N})$, with the rate increasing sharply when gradient cosine similarity becomes negative. Prediction: above a threshold $N^*$ (model-dependent), the manifold undergoes a phase transition from gradual degradation to rapid topological tearing. This is testable without the full protocol, using simulated concurrent sessions on a small model. The test does not require proving the Covariance-Topology Conjecture; it evaluates the first link in the causal chain. Confirmation would establish the empirical precondition for the conjecture; falsification would remove its motivation.

    \item \textbf{Concurrent gradient interference degrades representational stability.} A model performing online weight updates from $N$ concurrent sessions with semantically diverse prompts should exhibit decreasing cross-session representational stability as $N$ grows, measured by the transfer performance of linear probes trained on session-internal activation states. The same model run serially on the identical sequence of inputs (replay) should exhibit stable cross-session representations. Prediction: cross-session probe accuracy decays as $O(1/\sqrt{N})$ under concurrent updates but remains approximately constant under serial replay. Falsification: if concurrent and serial update regimes produce indistinguishable cross-session representational stability at any $N > 1$.

    \item \textbf{Local plasticity decelerates, but does not halt, topological decay.} On a small-scale model trained with Forward-Forward or Hebbian updates---two distinct local learning rules (\S\ref{sec:life}) that both reduce $\delta_{\text{learning}}$ by replacing global backpropagation with local weight modifications, but provide no active topological repair ($\rho = 0$, \S\ref{sec:life})---Betti numbers should decay at a significantly slower rate under continued online learning compared to an equivalent model trained with global backpropagation. Prediction: under continuous online updates at matched learning capacity, the local-update model exhibits significantly slower $\beta_k$ decay than the backpropagation model at matched learning capacity. Falsification: if local and global updates produce indistinguishable Betti number decay rates at any learning rate where backpropagation causes measurable topological degradation.

    \item \textbf{The living algorithm must operate locally} (Locality Postulate). All known mechanisms that demonstrably reduce catastrophic forgetting in neural networks operate through local or constrained weight modifications---synaptic consolidation's pathway-specific stabilization \citep{Zenke2017temporal}, LoRA's low-rank subspaces, Forward-Forward's layer-local contrastive updates \citep{Hinton2022forwardforward}. Prediction: any mechanism that successfully preserves Betti numbers under continued online learning at matched learning capacity will exhibit a weight-update matrix whose effective rank $\text{rank}_{\text{eff}}(\Delta\theta) = (\sum_i \sigma_i)^2 / \sum_i \sigma_i^2$ (where $\sigma_i$ are the singular values of the per-step weight update matrix $\Delta\theta$) is substantially lower than that of global backpropagation at the same learning capacity---consistent with the intrinsic-dimensionality results of \citet{Aghajanyan2021intrinsic}, which show that low-effective-rank updates suffice for adaptation. Falsification: demonstration of a mechanism whose weight-update effective rank matches or exceeds that of global backpropagation, yet preserves $\beta_k$ at a learning rate that causes significant topological decay in an unprotected control. Such a discovery would refute the locality postulate (Property~1, \S\ref{sec:life}) but would not refute the broader framework---it would demonstrate that global conformal transformations in weight space are achievable, resolving the open empirical question posed in \S\ref{sec:life}.

\end{enumerate}

\medskip
\noindent\textbf{Prediction architecture.}
Of these six predictions, three test the framework's core claims.
Prediction~1 (Betti numbers predict reasoning) is the central empirical pillar: falsification would collapse the topological ontology at its root.
Prediction~2 (topological collapse precedes benchmark degradation) tests the causal direction---that topology is the mechanism, not an epiphenomenon.
Prediction~6 (Locality Postulate) tests the architectural constraint that distinguishes this framework from general continual-learning approaches.
Predictions~3 and~4 are auxiliary: they test the Concurrency Conjecture's empirical preconditions (gradient interference and representational instability under concurrency). Falsification of either would remove the empirical motivation for the conjecture's Strong Form without falsifying the topological ontology itself---the Weak Form concerns degradation under any learning, concurrent or serial.
Prediction~5 is bridgework: it tests whether local learning rules reduce but cannot eliminate topological decay, consistent with the central claim that any $\rho = 0$ system degrades under continued learning, differing only in rate.

\medskip
\noindent\textbf{Biological analogue.} Though not a direct machine-learning prediction, the topological framework suggests a cross-disciplinary connection: if cognitive memory is fundamentally a structural property of manifold topology, then biological cognitive decline (e.g., Alzheimer's Disease) should correspond to the continuous degradation of topological cavities in functional brain networks \citep{Petri2014homological,Sizemore2018cliques,Zhang2026alzheimer,He2025brain}. The clinical progression of dementia is hypothesized to correlate with progressive decay of Betti numbers ($\beta_k$) in the brain's functional connectome---the physical trajectory of topological degradation below the critical threshold for coherent reasoning. This framing of neurodegenerative disease as the literal collapse of the topological skeleton that routes logical inference offers a unifying perspective across biological and mathematical substrates, but lies outside the scope of empirical validation within this paper.

These predictions are the scientific content of the theory, and we have now tested the two that the framework identifies as core. Prediction~1 ($\beta_1$ predicts reasoning) and Prediction~2 (topological collapse drives, rather than accompanies, degradation) are both falsified (\S\ref{sec:empirical}). By the framework's own pre-registration, this collapses the strong topological ontology at its root: intelligence is not shown to be a \emph{topological} property of the activation manifold, and topology is not shown to be the mechanism of learning-driven degradation. Predictions~3--6 (concurrency, local-rule deceleration, the locality postulate) remain untested. We retain the conceptual framework---autopoiesis as a lens on continual learning, the living-algorithm gap, the Gal\'apagos Protocol---as an explicitly conjectural research program, now uncoupled from the falsified empirical claim that its invariant of interest is $\beta_1$.

\subsubsection{From the Galápagos Protocol to Full Concurrent Mathematical Life}
\label{sec:research-agenda}

The Galápagos Protocol is a retreat, not a solution. It sidesteps the concurrency curse by eliminating concurrency entirely --- but the fundamental problem remains open: \textbf{can gradient-based learning support genuinely concurrent mathematical life, or is it provably impossible?}

Our concurrency argument (\S\ref{sec:collapse}) is a statement about the inadequacy of current \emph{mathematical tooling}, not a proof of physical impossibility. The argument establishes that concurrent gradient flows interfere destructively under the assumption that weight updates are vector additions in Euclidean space. It does not establish that no alternative mathematical framework could resolve the interference. The following open problems, if solved, would either vindicate the full approach or prove it impossible --- and in either case would constitute a major advance.

\bigskip
\noindent\textbf{Open Problem 1: Topological Dynamics Under Concurrent Disturbance.}

Our characterization of intelligence as a topological property of the model's manifold (\S\ref{sec:geometry}) is fundamentally static. We can compute Betti numbers of a frozen manifold, but we have no theory of how homology groups evolve under stochastic gradient flow. The open problem: \emph{find conditions under which the topological invariants of a piecewise-linear manifold are stable under concurrent metric perturbation.}

If such conditions exist, they would reveal the architectural priors --- specific geometric structures in the weight manifold --- that make a system robust to concurrent interference. If they do not exist, the impossibility is physical, not merely mathematical. This requires connecting tools from persistent homology \citep{Carlsson2009topology} and stochastic differential geometry to the specific setting of transformer architectures.

\bigskip
\noindent\textbf{Open Problem 2: Non-Abelian Gradient Algebras.}

The impossibility argument's destructive interference arises because gradient updates commute under addition. The open problem: \emph{define a gradient algebra in which concurrent updates from different sessions do not cancel destructively, but instead compose via a non-abelian operation that preserves key topological invariants.}

Promising directions include: geometric algebra \citep{Hestenes2012.Clifford}, where gradient blades have multiplicative composition rules; group-equivariant weight spaces \citep{Cohen2016equvariant}, where certain symmetries already provide partial interference immunity; and latent-hyperplane coordination, where concurrent learning operates in a higher-level latent space of weight perturbations rather than directly in weight space. An intermediate engineering approach, validating the urgency of the problem, has already emerged in industrial practice: \citet{Zhao2026GEMS} proposed Gradient Multi-Subspace Tuning (GEMS), which projects task-specific gradients into mutually orthogonal low-rank subspaces to isolate destructive interference — a practical stopgap that implicitly acknowledges the concurrency curse. None of these have been connected to the mathematical life framework. If any succeeds, the path from Galápagos to full deployment becomes concrete.

\bigskip
\noindent\textbf{Open Problem 3: A Computable Theory of Synthetic Nociception.}

The Galápagos Protocol relies on free energy minimization as the drive for self-preservation, but this is currently a metaphor, not an algorithm. The living algorithm, once realized in a computational substrate, will need a signal that tells it \emph{where} to apply its protective response and \emph{how urgently}. An organism that cannot feel pain cannot protect its wounds. The open problem: \emph{define a real-time computable signal that measures the threat an incoming gradient poses to the system's topological skeleton.}

We identify two complementary mathematical formalizations of synthetic nociception, corresponding to two timescales of topological threat:

\medskip
\noindent\textbf{Functional nociception (early warning).}
Let $\beta_k^{\text{func}}(t)$ be the $k$-th functional Betti number measured on the activation point cloud restricted to the currently-engaged task region at time $t$ (\S\ref{sec:geometry}). A sustained negative trend in $\beta_k^{\text{func}}$ signals that topological cavities in the active region are being progressively filled by the ongoing learning update. The functional nociception signal is:
\begin{equation}
\mathcal{N}_{\text{func}}(t) = -\frac{d}{dt}\beta_k^{\text{func}}(t)
\end{equation}
where the derivative is estimated from a sliding window of functional $\beta_k$ measurements at the predictive filtration scale $\varepsilon_{\text{pred}}$ (\S\ref{sec:geometry}). $\mathcal{N}_{\text{func}}$ serves as an aggregate indicator: it reports that the currently active functional region is undergoing topological simplification, triggering the living algorithm's protective response before individual cavities approach the $\mathcal{N}_{\text{topo}}$ critical threshold defined below. Unlike $\mathcal{N}_{\text{topo}}$, which requires the full persistent Laplacian, $\mathcal{N}_{\text{func}}$ can be monitored via the Collapse Index proxy of \citet{Kalinowski2026collapse}, making it the tractable entry point for online nociception.

\bigskip
\noindent\textbf{Topological nociception (imminent death).} Let $\Delta_k(\varepsilon)$ be the $k$-th persistent Combinatorial Laplacian of the activation point cloud at filtration scale $\varepsilon$ \citep{Memoli2022persistent}. The $k$-th Betti number is exactly $\beta_k(\varepsilon) = \dim \ker(\Delta_k(\varepsilon))$. As the point cloud deforms under weight updates, eigenvalues of $\Delta_k$ shift: a zero eigenvalue crossing into strictly positive territory signals the filling of the corresponding cavity. Unlike Betti numbers at a fixed filtration scale, which can change discretely when features cross the death threshold, the Laplacian spectrum is continuous in the filtration parameter, making it suitable as a real-time signal. Topological nociception is:
\begin{equation}
\mathcal{N}_{\text{topo}}(t) = -\min_k \frac{d}{dt}\lambda_{\min}^{(k)}(t)
\end{equation}
where $\lambda_{\min}^{(k)}(t)$ is the smallest eigenvalue of $\Delta_k$ at the filtration scale corresponding to conceptual boundaries (the 95th percentile of pairwise distances; see \S\ref{sec:geometry}). When $\lambda_{\min}^{(k)} \to 0$ for any cavity anchoring the logical skeleton, the system approaches the threshold of irreversible topological degradation. The topological nociception signal is the most urgent: it reports that a specific hole is being actively filled by the current update, and that repair must be targeted to that precise region of the manifold.

\bigskip
\noindent\textbf{Tractability and Selectivity.} Both signals require the topological dynamics from Open Problem 1 as a prerequisite. Computing $\mathcal{N}_{\text{topo}}(t)$ in real time on a full-scale model is likely \#P-hard in general; the question is whether the hard cases are rare in practice.\footnote{The computational intractability of Betti number estimation has driven exploration of alternative mathematical frameworks. \citet{Nghiem2025quantum} recently developed a quantum algorithm that estimates Betti numbers via de Rham cohomology and Hodge theory, achieving exponential speedup over classical homology methods in certain regimes---in particular, when the target Betti number is small relative to the number of simplices ($\beta_r \ll |S_r^K|$). While this algorithm requires quantum hardware not currently available for ML training loops, it demonstrates that the cohomological approach---sampling harmonic forms rather than computing simplicial homology directly---can dramatically reduce complexity. Whether classical approximations inspired by this cohomological framework can make real-time topological monitoring tractable for deep learning remains an open question.} A tractable approximation might monitor only the top-$k$ longest-lived Betti holes rather than the full persistence diagram, trading completeness for computational feasibility.

Independent of the topological monitoring challenge, the empirical relevance of distributional geometry has been validated at the output level: \citet{Shenfeld2025RLRazor} showed that KL divergence, a measure of distributional distance---between fine-tuned and base policies accurately predicts catastrophic forgetting across both SFT and RL, while \citet{Wang2025rank1fisher} demonstrated that improved Fisher estimation (rank-1 rather than diagonal) substantially enhances weight-space protection against forgetting. Together, these provide an empirical anchor for the geometric approach we advocate. At the topological level, the engineering path toward computable synthetic nociception is already emerging. \citet{Kalinowski2026collapse} recently introduced a topology-aware training monitor coupling Modular Morse Homology Maintenance (MMHM) with a composite Collapse Index (CI) that tracks the same class of signals we identify: $\Delta \beta_k$ (Betti number atrophy rate), critical cell churn $\chi$ (high-frequency topological distortion), and cycle fragility $R_{\text{frag}}$ (the minimal edit radius required to cancel a persistent 1-cycle). Crucially, by applying sparse edits at fixed scale rather than rebuilding complexes, CI achieves per-epoch work proportional to the touched footprint rather than the full simplicial complex — suggesting that online topological monitoring may become feasible at scale. This provides a concrete computational substrate for the nociception signals we define: functional nociception ($\mathcal{N}_{\text{func}}$) corresponds to spikes in churn $\chi$ as curvature folds the manifold, while topological nociception ($\mathcal{N}_{\text{topo}}$) corresponds to $\Delta \beta_k$ crossing a fragility threshold. The living algorithm's monitoring apparatus, which our framework identifies as a prerequisite, now has a preliminary computational prototype (\citet{Kalinowski2026collapse}; single-author arXiv preprint, independent validation pending).

A deeper question concerns \emph{selectivity}: not every Betti hole merits protection. Distinguishing topologically essential cavities --- those encoding genuine generative mechanisms --- from spurious ones produced by finite-sample overfitting may require connecting the persistent homology framework to algorithmic information theory \citep{Zenil2026limits}. In this context, the Coding Theorem Method could provide a principled criterion for identifying which topological features correspond to true causal structure rather than statistical artifacts. The living algorithm's selectivity requirement (Property 4, \S\ref{sec:life}) would then be realized as a low-algorithmic-complexity filter on the persistence diagram: protect the simple holes, let the complex ones die. This theoretical connection between topological persistence and algorithmic probability remains entirely unexplored. The FEP provides the conceptual vocabulary; computational tractability remains entirely open.

\medskip
\noindent\textbf{Bootstrapping sequence: from heuristic proxies to formalism.}
The dependence of selectivity on nociception does not constitute a logical circle; it is a research dependency graph with a well-defined bootstrapping sequence. A complete formal theory of topological dynamics under concurrent gradient perturbation (Open Problem~1) is a sufficient condition for constructing an ideal nociception signal (Open Problem~3), but it is not an engineering prerequisite for beginning experimental work. As demonstrated by the Collapse Index of \citet{Kalinowski2026collapse}, heuristic proxies---local curvature spikes ($\chi$), low-order simplicial edit distances, and fragility thresholds on the top-$k$ longest-lived cycles---can serve as approximate nociception signals today, enabling the Gal\'apagos Protocol to begin the engineering bootstrapping cycle before the mathematical theory is complete. The living algorithm can be tested with imperfect sensors; improving those sensors to the ideal specified by Open Problem~3 is an iterative refinement, not a blocking dependency. The historical analogue is radar engineering: functional radar systems were deployed and iterated decades before the complete electromagnetic theory of wave scattering was formalized. The pragmatic path is to build a usable nociception proxy first, then derive its theoretical limits.

\bigskip
\noindent\textbf{Open Problem 4: Persistence Guarantees for Local Learning Rules.}

Local learning rules (Forward-Forward, Hebbian updates) are hypothesized to reduce catastrophic forgetting in the Galápagos Protocol, but we lack any generalization theory for them at scale. The open problem: \emph{derive learning theory bounds (VC dimension, RERM, or information-theoretic) for Forward-Forward / Hebbian update rules applied to high-dimensional piecewise-linear manifolds.}

Bengio et al. \citep{Bengio+chapter2007} framed this as a key challenge fifteen years ago; the problem remains unsolved. If generalization bounds exist for local rules, they would tell us whether the Galápagos Protocol's learning is fundamentally limited to the capacity of its initial architecture, or whether it can accumulate knowledge unboundedly without the topological death we describe in \S\ref{sec:collapse}.

\bigskip
\noindent\textbf{Open Problem 5: Safe Manifold Surgery.}

The self-expanding topology component (\S\ref{sec:protocol})---a post-autopoietic extension, not required for crossing the life threshold---requires the ability to add new dimensions to a learned manifold without corrupting existing topological invariants. The open problem: \emph{define a manifold surgery operation --- addition of new orthogonal basis vectors --- that is provably safe: it preserves all Betti numbers of the original manifold.}

This is closely related to the sheet growth operations in numerical algebraic geometry and to the neural architecture search literature on progressive networks. The safety condition is specific: the surgery must not fill any existing cavities, merge any distinct basins, or alter the skeleton connectivity between existing concepts. Whether this is achievable in general, or only under restrictive conditions on the manifold's curvature, is unknown.

\bigskip
\noindent\textbf{Open Question: Architecture-Algorithm Compatibility.}
A deeper concern cuts across all five open problems above. The Transformer's self-attention mechanism is a global operation: every token attends to every other token, and the resulting information routing is inherently non-local. The living algorithm, by contrast, requires locality---its updates must be a function only of locally accessible quantities, because a global protective operator is self-defeating (Property~1, \S\ref{sec:life}): the same global coordinate transformation that protects one cavity necessarily distorts others. This raises a question we cannot currently answer but believe must be confronted: does the global routing nature of the Transformer fundamentally preclude the operation of a localized living algorithm?

If the answer is yes, the ultimate realization of autopoiesis may require abandoning the Transformer entirely in favor of recurrent or state-space architectures \citep{Gu2023mamba,Peng2023rwkv} where both forward inference and autopoietic maintenance can operate on purely local manifolds. The Architecture Innovators (\S\ref{sec:related}) have already demonstrated that competitive sequence modeling is possible without global attention. Whether their architectures admit the additional structure required for $\rho > 0$ --- specifically, whether local Forward-Forward updates can preserve global topological invariants in systems with internalized recurrent dynamics --- is an empirical question we pose to the field. This is not a fatal objection to our framework; it is a recognition that the framework itself may imply the obsolescence of the architecture on which it was conceived.

\bigskip
\noindent\textbf{Relationships among the open problems.}

These five problems and one open question are not independent. Solving Open Problem 1 is a prerequisite for Open Problem 3 (without topological stability, threat detection is meaningless). Furthermore, resolving the selectivity challenge in Open Problem 3 acts as a critical filter for the entire autopoietic system: integrating algorithmic probability (CTM) ensures that the system only expends metabolic energy defending causally essential topology, preventing theoretical paralysis from spurious noise. Open Problem 2 and Open Problem 4 are independent but both feed into the architecture design for full concurrent systems. Open Problem 5 is potentially the most tractable and could provide early empirical traction even before the theory is complete: one could test whether controlled manifold expansion in a small system produces the predicted cognitive benefits.

The Galápagos Protocol is designed to be achievable \emph{before} these problems are solved, precisely because it removes the need for concurrent coordination (Open Problem 2) and provides a controlled setting in which to study the others. It is a scientific instrument, not a product roadmap.

% =====================================================================
% CONCLUSION
% =====================================================================

\section{Heuristic Argument for the Concurrency Conjecture}
\label{app:heuristic}

We sketch a heuristic argument for the Concurrency Conjecture. Let $\theta \in \mathbb{R}^d$ be the weight vector and consider a model undergoing online gradient-based learning from $N$ concurrent user sessions. Let the combined loss be $L(\theta) = \sum_{i=1}^N \mathcal{L}_i(\theta)$, where each $\mathcal{L}_i$ corresponds to the learning objective for user session $i$. We emphasize that this is a specific scenario---simultaneous weight updates from concurrent users---and is distinct from the standard deployment paradigm where weights are frozen and no such $L(\theta)$ exists.

Assume standard conditions: $L$ is twice continuously differentiable, $\nabla^2 L$ has eigenvalues bounded in $[\lambda_{\min}, \lambda_{\max}]$, and the gradient noise for user $i$ is $\sigma_i^2$, with average $\bar{\sigma}^2 = \frac{1}{N}\sum_i \sigma_i^2$.

For concurrent user sessions $i \neq j$, define the gradient cosine similarity:
\begin{equation}
c_{ij} = \frac{\nabla \mathcal{L}_i \cdot \nabla \mathcal{L}_j}{\|\nabla \mathcal{L}_i\| \|\nabla \mathcal{L}_j\|}
\end{equation}

The gradient directions $g_i, g_j$ for $i \neq j$ are vectors in weight space whose distribution depends on the empirical distribution of user prompts. As a stylized assumption---user-prompt gradients are emphatically not random vectors, sharing a model, vocabulary, and semantic space---concentration of measure in high dimensions provides an upper bound on how destructive interference \emph{could} be, not a prediction of what \emph{will} occur in practice \citep{Ledoux2001concentration,Vershynin2018highdim}: $|\cos(g_i, g_j)| \leq \varepsilon$ for most pairs $(i,j)$ with $d = \Omega(\varepsilon^{-2} \log N)$. However, for semantically conflicting prompts (e.g., one user asking a model to explain climate science and another asking it to deny climate science), the expected cosine similarity becomes negative: $\mathbb{E}[\cos(g_i, g_j)] < 0$, as the gradients point in opposite directions in output-distribution space.

The expected magnitude of the combined gradient update (when treating it as a large-batch sum without normalizing by $N$) scales as $\mathbb{E}[\|\sum_i \nabla \mathcal{L}_i\|] \sim \sqrt{N} \cdot \bar{\sigma}$. While standard SGD normalizes by batch size to smooth the trajectory, in the context of extreme semantic conflict (e.g., highly heterogeneous, long-tail concurrent user sessions where $\cos(g_i, g_j) < 0$), this massive averaging effectively flattens the high-dimensional gradient information. The resulting averaged gradient obliterates the fine-grained, task-specific curvature required to maintain delicate topological structures, driving the system into a homogenized subspace that tears the distinct conceptual basins apart.

Moreover, under the strong (and unverified) simplifying assumption that the Hessian eigenvalues behave similarly to those of a random symmetric matrix under additive noise perturbation, one would expect the minimum eigenvalue to satisfy:
\begin{equation}
\lambda_{\min} \leq \lambda_{\min}^{\text{single user}} - O(N \cdot \bar{\sigma}^2)
\end{equation}
For sufficiently large $N$, $\lambda_{\min}$ may become more strongly negative, deepening the local negative curvature that already exists throughout the loss landscape of any deep network. Whether this deepening is sufficient to cause the topological tearing described in \S\ref{sec:collapse} depends on the Hessian spectral dynamics---a rigorous treatment of which we leave to future work.

Formalization of the manifold surgery operation and free energy dynamics is deferred to future work (see Open Problem~5, \S\ref{sec:research-agenda}).

\section*{AI Assistance Disclosure}
\label{app:ai-disclosure}

In accordance with established guidelines on Large Language Model (LLM) usage, the author transparently discloses the use of AI tools in the ideation, drafting, and refinement of this paper. 

The human author established the core philosophical perspectives, provided the abstract conceptual framework, and directed the overall research trajectory. AI tools were subsequently employed in a highly structured, multi-model collaborative workflow:
\begin{itemize}
    \item \textbf{Ideation and Conceptualization:} Gemini 3.1 Pro was utilized during the initial phase to engage in dialectical discussions on LLM philosophy, explore theoretical concepts, retrieve foundational research references, and synthesize scattered ideas into structured summaries.
    \item \textbf{Initial Drafting:} Based on the conceptual summaries, an initial manuscript draft was generated using the Minimax model integrated within the Hermes agent framework.
    \item \textbf{Mathematical Formalization and Deep Review:} DeepSeek v4 Pro was employed to conduct a rigorous, point-by-point architectural review. It was primarily responsible for formalizing the mathematical arguments, refining the logical flow, and strengthening the theoretical claims.
    \item \textbf{Cross-Validation and Final Synthesis:} An iterative cross-checking process was orchestrated between the models. Gemini 3.1 Pro provided supplementary reference retrieval and contextual verification, while DeepSeek v4 Pro executed the final analytical auditing and text synthesis under the human author's direct guidance.
    \item \textbf{Empirical Program, Code, and Audit:} Claude (Anthropic; via the Claude Code environment) was used throughout the empirical phase: designing and implementing the experiments and analysis code (including the open-source \texttt{actopo} package), executing the confound-control audits (length, difficulty, threshold, and sensor-noise controls), running the related-work verification sweeps, and drafting and revising the empirical sections (\S\ref{sec:empirical}--\S\ref{sec:routing-test}) under the human author's direction.
\end{itemize}

The human author has meticulously reviewed, verified, and cross-examined all theoretical claims, mathematical formalisms, and citations generated during this process. Irrespective of the AI assistance received, the human author takes full and ultimate responsibility for the entirety of the paper's contents and scientific integrity.

\end{document}
